\begin{textsample}{NEG Dim 3 – 2016\_c6 – Score -5.00 – 2016\_c6/enade\_2016\_med\_d7a7806f6b.txt}  \label{ex:f3_neg_035}
CAPÍTULO 3 – ANÁLISE TÉCNICA DA PROVA Este capítulo tem por objetivo apresentar o desempenho dos estudantes concluintes de Medicina no Enade/2016.
Para isso, foram calculadas as estatísticas básicas da prova em seu todo ( seção 3.1.1 ), bem como as estatísticas relacionadas aos Componentes de Formação Geral ( seção 3.1.2 ) e de Conhecimento Específico da Área ( seção 3.1.3 ).
Dadas as suas características, foram analisadas, em separado, as questões objetivas ( seção 3.2 ) e as questões discursivas ( seção 3.3 ).
Para as questões objetivas, foram disponibilizados os índices de facilidade e de discriminação ponto-bisserial, também em separado, para os Componentes de Formação Geral ( seção 3.2.1 ) e de Conhecimento Específico ( seção 3.2.2 ).
De cada componente, uma das questões foi escolhida para exemplificar a análise gráfica, relacionando as alternativas escolhidas pelos estudantes ( inclusive o gabarito ) com o número de acertos no componente.
O Anexo I apresenta a íntegra da análise gráfica para todas as questões objetivas.
Para cada uma das questões discursivas, os conteúdos dos tipos mais comuns de respostas dos estudantes são apresentados e comparados com o padrão de respostas esperado ( ver Anexo VIII com o padrão de respostas ).
Tomando como base as duas questões discursivas do Componente de Formação Geral, a seção 3.3.1.6 apresenta comentários sobre a correção das respostas com respeito à Língua Portuguesa.
Nas Tabelas 3.1 e 3.3 são apresentados os tamanhos da população inscrita, ausente, presente e percentual de faltosos, respectivamente, por Grande Região e por Categoria Administrativa e Organização Acadêmica da IES.
As demais tabelas apresentam as seguintes estatísticas das notas¹⁵ : média do desempenho na prova, erro padrão da média, desvio padrão, nota mínima, mediana e nota máxima para cada um de seus componentes.
Tais estatísticas contemplam o total de estudantes concluintes da Área de Medicina inscritos e presentes à prova do Enade/2016, tendo em vista agregações, ou por Grandes Regiões e o país como um todo, ou por Categoria Administrativa e Organização Acadêmica da IES.
Em relação aos gráficos de distribuição de notas, o intervalo considerado foi de 10 unidades, aberto à esquerda e fechado à direita, com exceção do primeiro intervalo, [ 0 ; 10 ], fechado em ambos os extremos.
Para os gráficos de distribuição das notas das ¹⁵ Uma definição dessas estatísticas pode ser encontrada no Glossário. questão em branco16 e nota zero. 3.1 - CONDIÇÃO DE PRESENÇA E ESTATÍSTICAS BÁSICAS DA PROVA Esta seção apresenta a condição de presença, além das estatísticas selecionadas e histogramas da nota geral ( 3.1.1 ) e de cada componente : Formação Geral ( 3.1.2 ) e Conhecimento Específico ( 3.1.3 ).
São também apresentadas estatísticas selecionadas de subpopulações caracterizadas por Grande Região, Categoria Administrativa e Organização Acadêmica. 3.1.1 - Condição de Presença e Estatísticas Básicas Gerais A Tabela 3.1 apresenta a Condição de Presença na prova, por Grande Região, dos estudantes concluintes de Medicina.
A população total de Inscritos foi de 16.180.
Destes, 15.866 estiveram Presentes, sendo 1,9 \% o índice de não comparecimento.
A Região de maior ausência foi a Centro-Oeste ( 5,0 \% ) e a de menor ausência foi a Sul ( 0,8 \% ).
Já a Tabela 3.2 apresenta as Estatísticas Básicas da prova, por Grande Região, dos estudantes concluintes de Medicina.
A Média das notas da prova como um todo ( nas seções seguintes serão analisados os Componentes de Formação Geral e de Conhecimento Específico ) foi 65,1, sendo que os alunos da região Norte obtiveram a Média mais baixa ( 61,8 ), e os da região Sul obtiveram a Média mais alta ( 67,4 ).
As demais médias foram : 65,8 na região Nordeste, 64,5 na região Sudeste e 66,2 na região Centro-Oeste.
O Desvio padrão para o Brasil como um todo foi 11,0, sendo o maior Desvio padrão encontrado na região Norte ( 11,4 ) e o menor, na região Sul ( 10,1 ), indicando uma dispersão um pouco menor das notas desta última região. 16 Nesse grupo estão incluídas também as respostas classificadas como nulas ou desconsideradas.
A região que obteve a maior nota Máxima foi a Sudeste ( 90,7 ), ao passo que a região que atingiu a menor nota Máxima foi a Nordeste ( 87,9 ).
A Mediana do Brasil como um todo foi 66,6, sendo a maior Mediana obtida na região Sul ( 68,9 ), e a menor obtida na Norte ( 62,7 ).
Na região Norte, a nota Mínima foi 16,5, na região Sul foi 5,8 e em todas as demais regiões a nota Mínima foi zero.
Considerando -se as notas segundo Grande Região, observa -se que existe diferença estatisticamente significativa ao nível de 95 \% entre a menor Média, obtida na região Norte ( 61,8 ), e as médias de todas as demais regiões17.
O comportamento das notas dos estudantes de todo o Brasil pode ser observado no Gráfico 3.1 que apresenta um histograma com a distribuição das mesmas.
Essa é uma distribuição unimodal com a moda no intervalo ( 60 ; 70 ], seguido pelo intervalo ( 70 ; 80 ]. 17 Todas as comparações deste capítulo utilizam os intervalos de 95 \%.
Os erros-padrão da média que possibilitam os testes estão disponíveis nas tabelas.
A Tabela 3.3 apresenta informações referentes à Condição de Presença desagregadas por Categoria Administrativa e por Organização Acadêmica.
Da população total de Inscritos, 10.101 são de IES Privadas e 6.079 de IES Públicas.
Em relação à Organização Acadêmica, a maior participação foi obtida por estudantes de Universidades ( 10.861 ), seguido de Faculdades ( 3.634 ) e Centros Universitários ( 1.685 ).
A Categoria Administrativa de maior ausência foi a Privada ( 2,0 \% ), e entre as Organizações Acadêmicas foi a dos Centros Universitários ( 4,2 \% ), que ficou acima da taxa de ausência nacional de 1,9 \%.
Já a Tabela 3.4 apresenta informações referentes à Média da nota final desagregadas por Categoria Administrativa e por Organização Acadêmica.
A Média das notas da prova como um todo foi 65,1.
Em relação à Categoria Administrativa, os estudantes das IES Públicas obtiveram Média mais alta ( 67,6 ), e os das IES Privadas obtiveram Média mais baixa ( 63,6 ), que a Média nacional.
Observa -se que existe diferença estatisticamente significativa entre as médias das notas das IES Públicas e Privadas.
A diferença entre as médias das regiões Sul e Norte ( 5,6 ), a maior e a menor Média, é superior à diferença entre IES Públicas e Privadas ( 4,0 ), caracterizando uma maior diversidade regional do que administrativa.
O Desvio padrão para as IES Públicas ( 10,0 ) foi inferior ao do Brasil como um todo ( 11,0 ), indicando uma dispersão um pouco menor das notas nessa Categoria Administrativa.
No tocante à Organização Acadêmica, as Universidades obtiveram Média mais alta que a nacional ( 65,5 e 65,1, respectivamente ).
As Faculdades obtiveram Média igual a nacional ( 65,1 ).
A média dos Centros universitários foi 62,4.
Constata -se que existe diferença estatisticamente significativa ao nível de 95 \% nas médias das notas dos estudantes dos Centros universitários, menor média, e as médias dos estudantes dos outros dois tipos de Organizações Acadêmicas. 3.1.2 - Estatísticas Básicas no Componente de Formação Geral A Tabela 3.5 apresenta as Estatísticas Básicas em relação ao componente da prova que avalia a Formação Geral dos estudantes concluintes.
Os alunos de todo o Brasil obtiveram desempenho médio de 60,0.
Quanto à variabilidade, o Desvio padrão das notas dos estudantes do Brasil, como um todo, foi 14,5.
A maior Média foi obtida na região Sul ( 61,8 ), e a menor, na região Norte ( 57,8 ).
As demais médias foram : 60,5 na região Nordeste, 59,5 na região Sudeste e 60,9 na região Centro-Oeste.
Já o maior Desvio padrão foi obtido nas regiões Norte e Sudeste ( 14,7 nas duas regiões ), e o menor, na região Sul ( 13,8 ).
Os demais desvios padrões foram : 14,1 na região Nordeste, e 14,0 da região Centro-Oeste.
A maior nota no Componente de Formação Geral da prova do Enade/2016 foi 97,4, obtida por, pelo menos, um aluno na região Sul.
A menor nota Máxima foi obtida na região Norte ( 94,2 ).
A Mediana do Brasil, como um todo, foi 61,6, sendo a menor Mediana encontrada na região Norte ( 59,0 ), e a maior, encontrada na região Sul ( 63,2 ).
A nota Mínima nessa parte foi diferente de zero apenas na região Norte ( 11,9 ).
Considerando -se as notas segundo Grande Região, observa -se que apenas não existe diferença estatisticamente significativa entre a maior Média das notas do Componente de Formação Geral, obtida na região Sul ( 61,8 ) e a Média obtida na região Centro-Oeste ( 60,9 ).
O Gráfico 3.2 propicia a avaliação do desempenho dos estudantes do Componente de Formação Geral, a partir do histograma da distribuição das notas correspondentes.
A distribuição é unimodal, com moda em ( 60 ; 70 ], assim como na prova, como um todo ( Gráfico 3.1 ).
Nota -se, ainda que, no Gráfico 3.2, as notas apresentam maior dispersão do que no Gráfico 3.1 ( distribuição das notas da prova ), confirmado pela comparação dos desvios padrões : 11,0 para a nota da prova como um todo e 14,5 para o Componente de Formação Geral.
Na Tabela 3.6 são apresentadas as informações referentes ao desempenho dos concluintes do Componente de Formação Geral, em diferentes agregações : Categoria Administrativa e Organização Acadêmica.
Observa -se que existe diferença estatisticamente significativa entre as médias dos tipos de Categoria Administrativa.
A maior Média foi obtida por estudantes de IES Públicas ( 62,7 ), com uma diferença estatisticamente significativa da obtida por estudantes de IES Privadas ( 58,3 ).
Nota -se que também há diferença estatisticamente significativa entre as médias dos Centro Universitários, a menor Média ( 57,6 ) e os outros dois tipos de Organização Acadêmica.
Já a diferença entre a maior Média, das Universidades ( 60,4 ), e a Média das Faculdades ( 60,0 ), não é estatisticamente significativa. 3.1.3 - Estatísticas Básicas do Componente de Conhecimento Específico A Tabela 3.7 apresenta as Estatísticas Básicas referentes ao Componente de Conhecimento Específico da área de Medicina.
A Média do desempenho dos alunos do Brasil, como um todo, foi 66,8.
A maior Média foi obtida na região Sul ( 69,3 ), e a menor, na região Norte ( 63,1 ).
As demais médias foram : 67,5 na região Nordeste, 66,1 na região Sudeste e 68,0 na região Centro-Oeste.
Quanto à variabilidade das notas, o Desvio padrão do Brasil, como um todo, foi 12,3, sendo o maior Desvio padrão observado na região Norte ( 12,7 ), e o menor, na região Sul ( 11,3 ).
Os demais desvios foram : 11,7 na região Nordeste, 12,6 na região Sudeste e 12,1 na região Centro-Oeste.
A Mediana das notas dos estudantes de todo o Brasil foi 68,3.
A maior Mediana ocorreu na região Sul ( 70,7 ), e a menor, na região Norte ( 64,2 ).
As demais medianas foram : 68,8 na região Nordeste, 67,8 na região Sudeste e 69,2 na região Centro-Oeste.
A nota Máxima do Brasil, como um todo, foi 96,7, sendo obtida por, pelo menos, um aluno na região Sudeste.
As demais notas máximas foram : 95,0 na região Norte, 95,2 na região Nordeste, 92,5 na região Sul e 94,0 na região Centro-Oeste.
A nota Mínima foi diferente de zero nas regiões Norte ( 4,8 ) e Sul ( 3,5 ).
Observa -se que existe diferença estatisticamente significativa da menor Média das notas do Componente de Conhecimento Específico da região Norte ( 63,1 ) em relação às demais regiões.
Assim como os Gráficos 3.1 e 3.2, o Gráfico 3.3, apresentado a seguir, permite uma avaliação do desempenho de concluintes de Medicina em relação ao Componente de Conhecimento Específico com um histograma da distribuição das notas correspondentes.
Dentre as três distribuições apresentadas, esta é a mais concentrada nas notas altas.
Esta também é uma distribuição unimodal, e o intervalo modal é o ( 70 ; 80 ], seguido de perto pelo intervalo ( 60, 70 ], intervalo modal das distribuições anteriormente apresentadas.
A Tabela 3.8 apresenta uma comparação dos resultados em relação à Categoria Administrativa e à Organização Acadêmica, agora levando em conta o desempenho de alunos do Componente de Conhecimento Específico da prova.
No que se refere à Organização Acadêmica, a maior Média foi das Universidades ( 67,2 ), vindo, a seguir, a das Faculdades ( 66,8 ) e, depois, a dos Centros Universitários ( 64,0 ).
O maior Desvio padrão, e acima do valor para o Brasil, como um todo, foi o dos Centros Universitários ( 12,7 ).
As Faculdades obtiveram as maiores notas Máxima ( 96,7 ).
As Universidades obtiveram nota Máxima igual a 95,2 e os Centros universitários obtiveram nota Máxima igual a 92,8.
As medianas foram : 69,0 nas Universidades, 65,3 nos Centros Universitários e 68,2 nas Faculdades.
A nota Mínima foi zero para todas as Organizações Acadêmicas.
Observa -se que existem diferenças estatisticamente significativas ao nível de 95 \% no Componente de Conhecimento Específico entre as notas dos Centros universitários e das outras Organizações Acadêmicas.
Quanto à Categoria Administrativa, observa -se um comportamento semelhante àquele da parte de Formação Geral e à prova como um todo, ou seja, existe diferença estatisticamente significativa entre as médias das IES Públicas ( 69,3 ) e IES Privadas ( 65,3 ).
Neste caso também, a maior Média foi obtida por alunos de IES Públicas de ensino. 3.2 - ANÁLISE DAS QUESTÕES OBJETIVAS Esta seção apresenta estatísticas selecionadas e histogramas das Questões Objetivas de Formação Geral ( 3.2.1 ).
São também apresentadas e comparadas as médias das subpopulações caracterizadas por Grande Região. 3.2.1 - Componente de Formação Geral A Tabela 3.9 apresenta as Estatísticas Básicas relativas às oito questões objetivas do componente da prova que abrange a Formação Geral dos estudantes.
A Média do Brasil foi 56,7.
A menor Média foi encontrada na região Norte ( 54,4 ), e a maior, na região Sul ( 59,6 ).
As demais médias foram : 57,4 na região Nordeste, 55,8 na região Sudeste e 58,0 na região Centro-Oeste.
O Desvio padrão do Brasil foi 18,4, sendo o maior Desvio padrão encontrado na região Sudeste ( 19,0 ), e o menor, na região Centro-Oeste ( 17,0 ).
Os demais desvios foram : 18,4 na região Norte, 17,5 na região Nordeste e na Sul.
As medianas do Brasil como um todo e de quatro das cinco regiões foram 62,5 — a única região com Mediana diferente foi a Norte, com 50,0.
As notas máximas ( 100,0 ) e as notas mínimas ( 0,0 ) foram iguais para todas as regiões.
A Tabela 3.10 apresenta o Índice de Facilidade e o Índice de Discriminação ( ponto-bisserial ) para cada uma das questões objetivas do Componente de Formação Geral.
Quanto ao Índice de Facilidade, foram usadas as seguintes cores para diferenciar o nível de dificuldade da questão : • Azul para as questões classificadas com índice Muito fácil ( \&gt ; =0,86 ), verde para as questões classificadas com índice Fácil ( 0,61 a 0,85 ), amarelo para as questões classificadas com Médio ( 0,41 a 0,60 ), vermelho para as questões classificadas com Difícil ( 0,16 a 0,40 ) e roxo para as questões classificadas com Muito difícil ( \&lt ; =0,15 ).
Já quanto ao Índice de Discriminação, foram usadas as seguintes cores para qualificar a questão : • As questões classificadas com índice Fraco receberam a cor vermelho ( \&lt ; =0,19 ), as classificadas com Médio receberam a cor amarelo ( 0,20 a 0,29 ), as classificadas com Bom receberam a cor verde ( 0,30 a 0,39 ) e as classificadas com Muito bom ( \&gt ; =0,40 ) receberam a cor azul.
As questões objetivas do Componente de Formação Geral, segundo o Índice de Facilidade, foram assim avaliadas : das oito questões, uma teve o Índice de Facilidade tido como Muito fácil, questão 6 com 91 \% de acertos.
Três questões foram classificadas como Fácil, por terem índices de acertos iguais a 0,61 ( questão 7 ), 0,79 ( questão 2 ) e 0,85 ( questão 3 ).
Uma questão foi considerada com índice de dificuldade Médio, situando -se no intervalo entre 0,41 e 0,60 do Índice de Facilidade, já que teve 56 \% de acertos, enquanto duas questões foram classificadas na categoria Difícil, situando -se no intervalo entre 0,16 e 0,40.
Por fim, uma questão apresentou exatamente 15 \% de acertos, razão pela qual seria classificada como Muito difícil ( \&lt ; =0,15 ).
Como já comentado, para análise das questões objetivas, relativas à Formação Geral, segundo o poder de discriminação, utilizou -se o Índice de Discriminação ( ponto- A Tabela 3.10 apresenta o Índice de Facilidade e o Índice de Discriminação ( ponto-bisserial ) para cada uma das questões objetivas do Componente de Formação Geral.
Quanto ao Índice de Facilidade, foram usadas as seguintes cores para diferenciar o nível de dificuldade da questão : • Azul para as questões classificadas com índice Muito fácil ( \&gt ; =0,86 ), verde para as questões classificadas com índice Fácil ( 0,61 a 0,85 ), amarelo para as questões classificadas com Médio ( 0,41 a 0,60 ), vermelho para as questões classificadas com Difícil ( 0,16 a 0,40 ) e roxo para as questões classificadas com Muito difícil ( \&lt ; =0,15 ).
Já quanto ao Índice de Discriminação, foram usadas as seguintes cores para qualificar a questão : • As questões classificadas com índice Fraco receberam a cor vermelho ( \&lt ; =0,19 ), as classificadas com Médio receberam a cor amarelo ( 0,20 a 0,29 ), as classificadas com Bom receberam a cor verde ( 0,30 a 0,39 ) e as classificadas com Muito bom ( \&gt ; =0,40 ) receberam a cor azul.
As questões objetivas do Componente de Formação Geral, segundo o Índice de Facilidade, foram assim avaliadas : das oito questões, uma teve o Índice de Facilidade tido como Muito fácil, questão 6 com 91 \% de acertos.
Três questões foram classificadas como Fácil, por terem índices de acertos iguais a 0,61 ( questão 7 ), 0,79 ( questão 2 ) e 0,85 ( questão 3 ).
Uma questão foi considerada com índice de dificuldade Médio, situando -se no intervalo entre 0,41 e 0,60 do Índice de Facilidade, já que teve 56 \% de acertos, enquanto duas questões foram classificadas na categoria Difícil, situando -se no intervalo entre 0,16 e 0,40.
Por fim, uma questão apresentou exatamente 15 \% de acertos, razão pela qual seria classificada como Muito difícil ( \&lt ; =0,15 ).
Como já comentado, para análise das questões objetivas, relativas à Formação Geral, segundo o poder de discriminação, utilizou -se o Índice de Discriminação ( ponto- O Índice de Facilidade variou de 0,15 a 0,91, e o de Discriminação, de 0,29 a 0,53.
As cinco questões com Índice de Discriminação Muito bom figuraram entre níveis de dificuldade Fácil, Médio e Difícil : três classificadas na categoria Fácil ( a questões 2, 3 e 7 ) do Índice de Facilidade, uma na categoria Médio ( questões 1 ) e uma na categoria Difícil ( questões 4 ).
Em particular, a questão 7 foi a que apresentou o maior poder discriminatório, com índice 0,53, e foi considerada Fácil em termos de facilidade, com uma proporção de 0,61 acertos.
O máximo de acertos foi alcançado pela questão 6 com um Índice de Facilidade de 0,91.
A questão de número 8 apresentou um Índice de Facilidade de 0,28, ou seja, pouco mais do que um quarto dos estudantes conseguiu resolvê -la, dentro do universo de participantes.
Seu Índice de Discriminação foi Bom ( 0,39 ).
Já a questão 4 obteve Índice de Discriminação Muito bom, 0,48, e seu Índice de Facilidade foi Difícil ( 0,40 ).
O Gráfico 3.4, para exemplificar, analisa o comportamento da questão de número 2 de Formação Geral.
Trata -se de uma questão considerada Fácil, com relação à facilidade e a que obteve o segundo maior Índice de Discriminação dessa parte da prova.
Nesse gráfico, cada uma das cinco curvas representa o percentual de respostas em determinada alternativa da questão, em função do número de acertos dos bisserial ).
Nessa análise, as questões foram assim avaliadas : cinco das oito questões apresentaram índice acima ou igual a 0,40 e, assim, foram classificadas com índice Muito bom para esse grupo de estudantes.
Duas questões tiveram Índice de Discriminação Bom, com valor entre 0,30 e 0,39, e uma questão foi classificada com índice de discriminação Médio ( entre 0,20 e 0,29 ).
Nenhuma questão teve nível Fraco de discriminação para esse grupo de estudantes.
O Índice de Facilidade variou de 0,15 a 0,91, e o de Discriminação, de 0,29 a 0,53.
As cinco questões com Índice de Discriminação Muito bom figuraram entre níveis de dificuldade Fácil, Médio e Difícil : três classificadas na categoria Fácil ( a questões 2, 3 e 7 ) do Índice de Facilidade, uma na categoria Médio ( questões 1 ) e uma na categoria Difícil ( questões 4 ).
Em particular, a questão 7 foi a que apresentou o maior poder discriminatório, com índice 0,53, e foi considerada Fácil em termos de facilidade, com uma proporção de 0,61 acertos.
O máximo de acertos foi alcançado pela questão 6 com um Índice de Facilidade de 0,91.
A questão de número 8 apresentou um Índice de Facilidade de 0,28, ou seja, pouco mais do que um quarto dos estudantes conseguiu resolvê -la, dentro do universo de participantes.
Seu Índice de Discriminação foi Bom ( 0,39 ).
Já a questão 4 obteve Índice de Discriminação Muito bom, 0,48, e seu Índice de Facilidade foi Difícil ( 0,40 ).
O Gráfico 3.4, para exemplificar, analisa o comportamento da questão de número 2 de Formação Geral.
Trata -se de uma questão considerada Fácil, com relação à facilidade e a que obteve o segundo maior Índice de Discriminação dessa parte da prova.
Nesse gráfico, cada uma das cinco curvas representa o percentual de respostas em determinada alternativa da questão, em função do número de acertos dos estudantes nessa parte da prova ( Formação Geral/Múltipla Escolha ), antes de possíveis eliminações pelo critério do ponto-bisserial.
A curva em azul corresponde à alternativa A, a alternativa correta para esta questão.
Observa -se que entre os estudantes com menor número de acertos, nessa parte do exame, a situação mais frequente foi a escolha da alternativa B ( em verde ), incorreta.
Por exemplo, entre os estudantes que acertaram três questões, 27,9 \% escolheram a alternativa B ( em verde ), 11 \% escolheram a alternativa E ( vermelho ), 2 \% escolheram C ( preto ) e 0,8 \% a D ( roxo ).
Entre os que acertaram três das questões de múltipla escolha de formação geral, 0 \% deixou a questão em branco e 0,2 \% marcou mais de uma alternativa, invalidando a questão.
À medida em que o número de acertos aumenta, indicando desempenho melhor nessa parte da prova, aumenta concomitantemente a proporção de estudantes que selecionaram a alternativa correta A, atingindo 100 \% para os estudantes com 8 acertos.
Essa análise permite verificar como a questão discriminou os grupos de desempenho, justificando o alto índice obtido na questão.
Cumpre notar que não é possível inferir desse gráfico nem o índice de facilidade ( que seria uma média da proporção ponderada pela quantidade de alunos com cada uma das notas ), nem o índice de discriminação ponto-bisserial, por razão equivalente.
No caso extremo, no qual a grande concentração dos acertos dos alunos fosse abaixo de 4, o índice de facilidade seria obrigatoriamente abaixo de 40 \% ( neste exemplo ).
Caso a concentração fosse em 6 acertos ou mais, o índice seria obrigatoriamente acima de 60 \%.
Os gráficos relativos às demais questões de Formação Geral constam do Anexo I. 3.2.2 - Componente de Conhecimento Específico A Tabela 3.11 apresenta as Estatísticas Básicas em relação às questões objetivas do Componente de Conhecimento Específico da prova de Medicina por Grande Região.
A Média do Brasil desse componente foi de 72,9.
A menor Média foi observada na região Norte ( 68,9 ), e a maior, na região Sul ( 75,8 ).
O Desvio padrão de todo o Brasil foi 13,3, sendo o menor Desvio padrão encontrado na região Sul ( 12,3 ), e o maior, na região Norte ( 13,7 ).
A Mediana de todo o Brasil foi 72,7, a mesma encontrada na região Sudeste.
Na região norte, a Mediana foi 68,2, a menor delas, e 77,3 nas demais regiões, Nordeste, Sudeste e Centro-Oeste.
A nota Máxima da prova foi 100,0, obtida nas questões objetivas do Componente de Conhecimento Específico, por, pelo menos, um aluno de todas as regiões.
Em todas as regiões a nota Mínima foi zero.
A Tabela 3.12 apresenta os Índices de Facilidade e Discriminação ( ponto-bisserial ) das questões objetivas do Componente de Conhecimento Específico para os estudantes de Medicina.
Para facilitar a diferenciação das questões, usaremos as mesmas cores da Tabela 3.10 para as diferentes classificações dos Índices de Facilidade e de Discriminação.
A partir do Índice de Facilidade obtido, pode -se concluir que apenas quatro das questões objetivas da prova foram considerados, pelo menos, Difícil : das 27 questões válidas, três foram classificadas como Difícil e uma como Muito difícil. 14 questões foram classificadas como Fácil, cinco consideradas como Médio e quatro como Muito Fácil.
Já quanto aos índices de discriminação das questões objetivas do Componente de Conhecimento Específico da prova, tem -se como resultado a seguinte classificação : 11 das 27 questões foram consideradas como boas, enquanto três delas tiveram Índice de Discriminação Muito bom.
Assim, a maioria das questões – 14 em 27 – os Índices de Discriminação foram Bom ou Muito bom.
Dentre as demais, oito delas foram classificadas como Médio e outras cinco como Fraco, sendo 13, por conseguinte, a quantidade de questões nos dois patamares mais baixos de discriminação.
Constata -se, assim, que a prova – no que se refere ao Componente de Conhecimento Específico – possuía boa capacidade de discriminar entre aqueles que dominam ou não o conteúdo.
Dentre as questões que alcançaram os maiores índices de discriminação, três delas, as de números 15, 16 e 27, foram classificadas com Índice de Discriminação Muito bom, situando -se no intervalo de 0,40 a 0,50 do índice.
As três também foram classificadas como Fácil, quanto ao Índice de Facilidade.
A questão de número 26 foi a mais difícil dentre as 27 questões específicas, com baixo Índice de Facilidade, apenas 19,0 \% de acertos.
Essa questão apresentou poder discriminatório igualmente baixo, 0,09, o que comprova ter sido esta questão a mais difícil para os estudantes.
Destaca- A Tabela 3.12 apresenta os Índices de Facilidade e Discriminação ( ponto-bisserial ) das questões objetivas do Componente de Conhecimento Específico para os estudantes de Medicina.
Para facilitar a diferenciação das questões, usaremos as mesmas cores da Tabela 3.10 para as diferentes classificações dos Índices de Facilidade e de Discriminação.
A partir do Índice de Facilidade obtido, pode -se concluir que apenas quatro das questões objetivas da prova foram considerados, pelo menos, Difícil : das 27 questões válidas, três foram classificadas como Difícil e uma como Muito difícil. 14 questões foram classificadas como Fácil, cinco consideradas como Médio e quatro como Muito Fácil.
Já quanto aos índices de discriminação das questões objetivas do Componente de Conhecimento Específico da prova, tem -se como resultado a seguinte classificação : 11 das 27 questões foram consideradas como boas, enquanto três delas tiveram Índice de Discriminação Muito bom.
Assim, a maioria das questões – 14 em 27 – os Índices de Discriminação foram Bom ou Muito bom.
Dentre as demais, oito delas foram classificadas como Médio e outras cinco como Fraco, sendo 13, por conseguinte, a quantidade de questões nos dois patamares mais baixos de discriminação.
Constata -se, assim, que a prova – no que se refere ao Componente de Conhecimento Específico – possuía boa capacidade de discriminar entre aqueles que dominam ou não o conteúdo.
Dentre as questões que alcançaram os maiores índices de discriminação, três delas, as de números 15, 16 e 27, foram classificadas com Índice de Discriminação Muito bom, situando -se no intervalo de 0,40 a 0,50 do índice.
As três também foram classificadas como Fácil, quanto ao Índice de Facilidade.
A questão de número 26 foi a mais difícil dentre as 27 questões específicas, com baixo Índice de Facilidade, apenas 19,0 \% de acertos.
Essa questão apresentou poder discriminatório igualmente baixo, 0,09, o que comprova ter sido esta questão a mais difícil para os estudantes.
Destaca- Destaca -se, também, a questão 21, com Índice de Facilidade 0,19, o que, em termos percentuais, corresponde a 19,0 \% de estudantes que responderam acertadamente.
Já 0,16 foi o seu de Índice de Discriminação.
Tais questões foram, portanto, pelo critério ponto-bisserial, consideradas inadequadas.
Por isso, as questões 21 e 26, além das questões 13, 17 e 35 foram eliminadas do cômputo da nota final.
A título de exemplo das análises do comportamento das questões objetivas, o Gráfico 3.5 analisa a questão 27 do Componente de Conhecimento Específico.
Esta questão foi considerada, pelas respostas dos estudantes avaliados na prova, como uma questão com nível de facilidade Fácil, 0,74, ou seja, 74,0 \% dos estudantes assinalaram acertadamente a opção B, correspondente ao gabarito.
Seu Índice de Discriminação foi igual a 0,43, classificado como Muito bom, também sendo esta questão a que apresentou o maior valor.
Nesse gráfico, cada uma das cinco curvas representa o percentual de respostas em determinada alternativa da questão 27, em função do número de acertos dos estudantes nessa parte da prova, antes de possíveis eliminações de questões pelo critério do ponto-bisserial.
A alternativa correta, B, representada no gráfico pela curva em verde, foi escolhida em maiores proporções pelos alunos com desempenho melhor nessa parte da prova.
Já as alternativas incorretas, também denominadas distratores, foram selecionadas, principalmente, por aqueles com notas mais baixas.
Neste caso também a soma não é sempre 100 \% por causa das questões não respondidas ou com mais de uma opção marcada.
Aqueles com nota zero, na sua quase totalidade, deixaram esta questão em branco ou marcaram mais de uma alternativa, comportamento considerado inválido.
A proporção de alunos que selecionou a resposta correta B aumenta gradativamente, chegando a atingir 100 \% para 24 acertos, enquanto a proporção dos que escolheram alternativas incorretas decai, a partir de oito acertos, como função do número de acertos nessa parte da prova.
Os gráficos relativos às demais questões do Conhecimento Específico constam do Anexo I. 3.3 - ANÁLISE DAS QUESTÕES DISCURSIVAS Esta seção apresenta estatísticas selecionadas e histogramas das Questões Discursivas de Formação Geral ( 3.3.1 ) e Conhecimento Específico ( 3.3.2 ).
São também apresentadas e comparadas as médias de subpopulações caracterizadas por Grande Região. 3.3.1 - Componente de Formação Geral As análises dos resultados de desempenho dos estudantes de Medicina nas duas questões discursivas relativas à Formação Geral encontram -se na Tabela 3.13 e no Gráfico 3.6.
Na Tabela 3.13, observa -se que a nota Média nesse conjunto de questões foi acima da obtida nas objetivas.
Os estudantes, de todo o Brasil, obtiveram, em Formação Geral, Média 56,7 nas questões objetivas e 64,9 nas questões discursivas.
Pode -se notar também que o Desvio padrão nesse conjunto de questões foi um pouco maior do obtido nas objetivas, 18,4 nas questões objetivas e 19,1 nas questões discursivas.
A maior Média foi obtida na região Centro-Oeste ( 65,3 ), e a menor, nas regiões Norte onde a Média foi 62,9.
A Mediana de todo o Brasil, neste componente, foi 68,5, a mesma que as obtidas nas regiões Sudeste e Sul.
Na região Norte, a Mediana foi menor, 67,5.
Nas regiões Nordeste e Centro-Oeste foi maior que o valor para o Brasil : 69,0 e 69,5, respectivamente.
A nota Máxima do Brasil ( 98,5 ) foi a mesma que a da região Sudeste.
A menor nota Máxima foi a da região Norte ( 96,5 ), e as demais foram : 97,0 na região Sul e 98,0 nas regiões Nordeste e Centro-Oeste.
A nota Mínima ( 0,0 ) foi a mesma em todas as regiões do Brasil.
O Gráfico 3.6 representa a distribuição das notas nas questões discursivas do Componente de Formação Geral.
A moda dessa distribuição ocorre no intervalo ( 70 ; 80 ].
Destaca -se também o intervalo [ 60 ; 70 ] com distribuição maior do que 25 \% do total de notas.
No intervalo [ 0 ; 10 ], com menos do que 5 \% das notas, incluem -se, além da nota zero, a frequência de alunos que deixaram esse tipo de questão em branco.
Na sequência, os resultados verificados para cada uma das questões discursivas de Formação Geral serão apresentados, estabelecendo -se relações com os conteúdos abordados em cada uma delas.
Os comentários da Banca de docentes corretores, a respeito do observado na correção das respostas dos estudantes, suas impressões e conclusões serão apresentados junto à análise de cada questão.
Cumpre esclarecer que, tendo em vista que as questões discursivas de Formação Geral são padronizadas, ou seja, constam de todas as provas, os comentários da Banca são os mesmos para todas as carreiras acadêmicas, sendo direcionados a todos os estudantes que participaram do Enade/2016.
A seguir, serão analisados os desempenhos da Área de Medicina nas duas questões discursivas de Formação Geral do Enade/2016, comparando -se os resultados obtidos com comentários para cada questão. 3.3.1.1 - Análise de Conteúdo da Questão Discursiva 1 do Componente de Formação Geral Os dados de Medicina, obtidos a partir das respostas à questão 1, encontram -se na Tabela 3.14 e no Gráfico 3.7.
Nessa questão – de melhor desempenho dentre as duas de Formação Geral – os alunos, de todo o Brasil, tiveram Média 74,6.
A maior Média para a questão 1 foi obtida nas regiões Sul e Centro-Oeste ( 74,8 ), e a menor, na região Norte ( 72,5 ).
Quanto à variabilidade das notas, o Desvio padrão de todo o Brasil foi 25,6.
O menor Desvio padrão foi obtido na região Sul ( 24,9 ), e o maior Desvio padrão foi obtido na região Norte ( 27,4 ).
As medianas de todas as regiões e do Brasil como um todo, foram iguais ( 80,0 ).
As notas máximas e mínimas da questão discursiva 1 foram, sem exceção, as mesmas para todas as regiões do Brasil, respectivamente, 100,0 e 0,0.
O Gráfico 3.7 mostra a distribuição das notas na questão discursiva 1 do Componente de Formação Geral.
Observa -se que a maior frequência corresponde aos alunos que obtiveram nota dentro do intervalo ( 70 ; 80 ], correspondendo à moda da distribuição.
Destaca -se, também, o intervalo ( 90 ; 100 ] com cerca de 23 \% das notas obtidas.
Os alunos que deixaram a questão em branco foram aproximadamente 5 \% do total. 3.3.1.2 - Comentários sobre a correção de Conteúdo das respostas à Questão Discursiva 1 A questão discursiva 1 alinha -se aos conteúdos programáticos previstos nas diretrizes curriculares, quanto às habilidades e às competências para a atualização permanente e aos conhecimentos sobre a realidade brasileira, mundial, tangenciando outras áreas do conhecimento.
O enunciado é claro e indica um tema de extrema relevância, qual seja, as migrações humanas e todo o complexo de problemas que esse movimento causa, quer na origem, quer no seu ponto de destino.
A questão é extremamente atual, maximizada pelos acompanhamentos midiáticos que espetacularizam cenas dramáticas do processo de fuga de massas humanas das guerras, como a morte de imigrantes.
Além do aspecto internacional, deve se ter em conta que o Brasil sofreu vários influxos migratórios, desde o forçado, pela transferência de escravos da África, até o incentivado, dos italianos, alemães e outros europeus para trabalhar na agricultura.
Isso sem considerar a própria colonização portuguesa originária e as migrações posteriores dos mesmos portugueses.
Posteriormente, o Brasil tornou -se um país de emigrantes que foram \textbf{aportar} a outros locais na busca de melhor situação econômica.
Muitos migraram para a América do Norte e para a Europa, principalmente.
Mais recentemente, por força do crescimento econômico pátrio, novas levas de imigrantes desembarcaram no Brasil, muitos deles provenientes do Haiti, embora tenhamos recebidos vizinhos latinos, como bolivianos, paraguaios, chilenos, argentinos, venezuelanos, peruanos, colombianos que também vieram em busca de uma vida melhor.
Como se verifica, a temática indicada no enunciado pertence à realidade conhecida amplamente pelos estudantes, pois, também divulgada por todos os meios de comunicação.
Trata -se de tema geral que, normalmente, não se inclui nas grades curriculares superiores, mas é de conhecimento geral e relacionado aos direitos humanos, este sim, ponto de integração às diretrizes curriculares.
A linguagem utilizada foi clara, indicando a existência dos fatos vinculados aos movimentos migratórios.
Por outro lado, o texto demanda uma reflexão sobre a realidade dos indivíduos que são compelidos ao afastamento dos seus lares por motivos externos, como guerras, conflitos civis, religiosos, perseguições das mais variadas, questões econômicas.
Como se trata de texto contendo uma ampla gama de possibilidades, é natural que ocorram interpretações diversas das esperadas.
Do mesmo modo, a contemporaneidade do tema encaminha o enunciado para ser qualificado como questão de média dificuldade, em comparação com as demais questões apresentadas.
O padrão de resposta indicou possíveis interpretações do texto.
No que concerne ao grau de profundidade, mostrou a necessidade de conhecimento de vários itens pertinentes à migração e que são refletidos nas várias possibilidades de resposta previstas pelo padrão.
Pela análise das provas realizadas, constatou -se que muitos foram os que identificaram os itens indicados no padrão de resposta.
Essa circunstância poderia ser justificada, tendo em vista que o tema vem sendo bastante difundido na mídia social, gerando elementos naturais para a resposta, bem como da história do país, que, intimamente, vincula -se ao processo migratório, quer ativo, remetendo pessoas para outros países, quer passivo, na recepção de pessoas.
Essa característica também permite responder aos termos postos quanto aos Estados que recebem os migrantes.
Assim, a indicação dos aspectos positivos e negativos também é bastante difundida.
As respostas medianas e boas apreenderam o sentido do enunciado, e indicaram os motivos expressos no padrão de respostas, como os ensejadores das imigrações, notadamente, para a Europa, que seria o foco do texto.
Por outro lado, também identificaram os problemas dos imigrantes, no mesmo padrão indicado para as respostas desejadas.
Alguns concluintes procuraram exemplos de comparação na realidade brasileira, como os clássicos imigrantes italianos, alemães e de outras nacionalidades, bem como os atuais haitianos e demais povos da América Latina.
Uma minoria não apreendeu o sentido do enunciado e, somente, apresentou elementos da realidade nacional.
Esses foram avaliados no conjunto de provas fracas.
As respostas demonstraram que o concluinte está atualizado, em regra, com o tema da imigração, indicando os fundamentos dessa movimentação populacional que, nos últimos tempos, tem se agravado pela incidência de conflitos armados em vários países, ensejadores de reflexos na segurança dos habitantes desses locais, bem como o surgimento de focos de doenças e ausência dos serviços públicos básicos, além da fome, decorrente da impossibilidade de produzir alimentos ou mesmo comerciá -los.
Do ponto de vista dos Estados que recebem o fluxo de imigrantes, também restou clara a percepção das dificuldades vinculadas à ausência de estrutura quanto aos serviços públicos ( saúde, educação, segurança, moradia, alimentação ) para receber um número expressivo de imigrantes.
Em época de recessão econômica, a presença de estrangeiros acarreta a competição pela busca de empregos, gera movimentos xenofóbicos, bem como choques culturais e políticos.
Deve -se ter em mira que esses movimentos têm um viés de insegurança, pelos países receptores, diante da crise causada por movimentos considerados terroristas, embora a imensa maioria não tenha qualquer vínculo com esses movimentos.
Os recentes atentados na Europa indicam que os seus autores têm origem em comunidades de imigrantes, embora muitos já tenham nascido nos países da comunidade europeia.
Essas características negativas, no entanto, paradoxalmente, convivem com aspectos positivos, o que encaminha a questão para a necessidade dos jovens imigrantes, que ocuparão postos de trabalho não desejados pelos habitantes locais, bem como, compensarão os baixos índices de crescimento populacional.
Existe outro reflexo positivo que está relacionado ao sistema previdenciário, que terá novos contribuintes nos trabalhadores estrangeiros que se integram ao sistema.
Esses pontos foram abordados e indicam que existe uma consciência dos concluintes dos cursos universitários quanto aos principais itens relacionados ao tema que foi posto para análise. 3.3.1.3 - Análise de Conteúdo da Questão Discursiva 2 do Componente de Formação Geral A Tabela 3.15 mostra que o desempenho médio dos estudantes na questão discursiva 2 ( média 53,5 ) foi inferior ao obtido na questão discursiva 1 ( média 74,6 ).
A região Centro-Oeste foi aquela cuja Média, também nessa questão, foi maior ( 53,8 ), e a de menor Média foi a região Norte ( 51,8 ).
Quanto à variabilidade das notas, o Desvio padrão de todo o Brasil foi 23,7, inferior ao obtido na questão discursiva 1 ( 25,6 ).
O maior desvio nessa questão foi obtido na região Centro-Oeste ( 24,6 ), enquanto o menor foi obtido na região Sul ( 23,1 ).
A Mediana de todo o Brasil foi 60,0, a mesma de todas as regiões.
As notas máximas ( 100,0 ) e as notas mínimas ( 0,0 ) também foram as mesmas em todas as regiões, sem exceção.
O Gráfico 3.8 mostra a distribuição das notas na questão discursiva 2 do Componente de Formação Geral.
Observa -se que a maior frequência corresponde aos alunos que obtiveram nota dentro do intervalo ( 60 ; 70 ]. 5 \% dos alunos deixaram a questão em branco.
Nota -se, ainda, que as notas ficaram mais concentradas em comparação à questão discursiva de número 1.
Isto pode ser constatado também pela comparação do desvio padrão das notas da questão discursiva 2 ( 23,7 ) e da questão discursiva 1 ( 25,6 ). 3.3.1.4 - Comentários sobre a correção de Conteúdo das respostas à Questão Discursiva 2 A questão discursiva 2 tem seu conteúdo integrado às Diretrizes Curriculares, uma vez que a proteção aos direitos humanos e às relações de gênero constituem matrizes fundamentais para conhecimento dos profissionais de nível superior.
O enunciado é claro e indica um tema de extrema relevância, qual seja, a análise da questão da mulher submetida a injustas agressões e os seus reflexos pessoais, sociais e familiares.
Também enfrenta a busca pela maior proteção no concernente às agressões, com a instituição de um modelo legal mais repressivo e que foi gerado por uma situação especifica de ataque à mulher, com consequências dramáticas para a vítima, Maria da Penha, cuja lei levou seu nome.
Deve ser recordado que os conflitos vinculados às relações familiares sempre foram circunscritos ao seu âmbito e, quando dele extravasavam, por força de agressões contundentes à integridade física de um dos partícipes da relação, em geral a mulher, raramente redundavam em punições relevantes, sendo comum o pagamento em cestas básicas de alimentos, como sanção.
Após a edição da lei Maria da Penha ( Lei no. 11.340, sancionada em 7 de agosto de 2006 ), além das sanções sofrerem modificações, houve a criação de modelos de proteção e atendimento especial para os casos de violência doméstica, procurando proteger as mulheres agredidas antes de danos maiores, muitas vezes irreversíveis, incluindo aí o feminicídio, também incluído no Código Penal, como circunstância qualificadora do crime de \textbf{homicídio}, e incluí -lo no rol dos crimes hediondos.
A lei Maria da Penha busca realizar, como indicado no seu pórtico, a Convenção sobre a Eliminação de Todas as Formas de Discriminação contra as Mulheres e a Convenção Interamericana para Prevenir, Punir e Erradicar a Violência contra a Mulher.
No seu texto, várias são as formas de violência doméstica e familiar contra a mulher, identificadas como sendo, entre outras : “ a violência física, entendida como qualquer conduta que ofenda sua integridade ou saúde corporal ; a violência psicológica, entendida como qualquer conduta que lhe cause dano emocional e diminuição da \textbf{autoestima} ou que lhe prejudique e perturbe o pleno desenvolvimento ou que vise a degradar ou controlar suas ações, comportamentos, crenças e decisões, mediante ameaça, constrangimento, \textbf{humilhação}, manipulação, isolamento, vigilância constante, perseguição contumaz, insulto, chantagem, ridicularização, exploração e limitação do direito de ir e vir ou qualquer outro meio que lhe cause prejuízo à saúde psicológica e à autodeterminação ; a violência sexual, entendida como qualquer conduta que a constranja a presenciar, a manter ou a participar de relação sexual não desejada, mediante intimidação, ameaça, coação ou uso da força ; que a induza a comercializar ou a utilizar, de qualquer modo, a sua sexualidade, que a impeça de usar qualquer método contraceptivo ou que a force ao matrimônio, à gravidez, ao aborto ou à prostituição, mediante coação, chantagem, suborno ou manipulação ; ou que limite ou anule o exercício de seus direitos sexuais e reprodutivos ; a violência patrimonial, entendida como qualquer conduta que configure retenção, subtração, destruição parcial ou total de seus objetos, instrumentos de trabalho, documentos pessoais, bens, valores e direitos ou recursos econômicos, incluindo os destinados a satisfazer suas necessidades ; a violência moral, entendida como qualquer conduta que configure calúnia, difamação ou injúria. ” ( Conforme texto extraído do artigo sétimo da referida lei ).
Também se buscam medidas de integração entre os vários órgãos públicos, abrangendo repressão, proteção, conscientização, bem como utilização de equipes multidisciplinares, tendo em vista os graves danos, de variados matizes, que redundam das práticas agressivas.
Nessa linha, também existe na lei a previsão de locais de atendimento específico ( centros de atendimento integral e multidisciplinar para mulheres e respectivos dependentes em situação de violência doméstica e familiar ; casas-abrigo para mulheres e respectivos dependentes menores em situação de violência doméstica e familiar ; delegacias, núcleos de defensoria pública, serviços de saúde e centros de perícia médico-legal especializados no atendimento à mulher em situação de violência doméstica e familiar ; programas e campanhas de enfrentamento da violência doméstica e familiar ; centros de educação e de reabilitação para os agressores ).
O Poder Judiciário, também, deve especializar o seu atendimento para aplicar as medidas de restrição ao agressor de forma mais eficiente.
A linguagem utilizada na questão foi clara, indicando a grave violação dos direitos das mulheres injustamente agredidas e a luta pela busca de sanções adequadas aos agressores e de proteção pessoal e para os de sua família.
Ademais, o tema não está circunscrito ao Brasil, mas o país possui um número muito expressivo de atos que compõem esse cenário.
Por outro lado, o texto demanda uma reflexão sobre a realidade, e abarca diversos segmentos da sociedade, pois as situações de agressão não se restringem às camadas mais populares.
Como se trata de texto contendo uma ampla gama de possibilidades, é natural que ocorram interpretações diversas das esperadas.
O padrão de resposta identifica a projeção do elaborador na resposta a ser apresentada pelo estudante.
Está em linguagem correta e indica possíveis interpretações do texto.
Suas opções foram bastante amplas, abrangendo a maior parte das respostas encontradas.
No que concerne ao grau de profundidade, o padrão de resposta indicou a necessidade de conhecimento de lei especial que regula a violência doméstica, o que exige ou qualificação da área de Direito ou que os currículos dos cursos dos participantes desse ciclo do Enade tenham oferecido alguma disciplina introdutória da área jurídica.
Ao analisar as provas realizadas, constatamos que, na primeira parte, houve maior facilidade em responder aos itens indicados no padrão de resposta, pois atinentes, não só à formação profissional, mas, também, ao conhecimento da realidade, pois, infelizmente, a violência doméstica está disseminada na sociedade e nos vários pontos do país.
Mas, verificou -se que as respostas vieram, no mais das vezes, desorganizadas, dificultando a correção.
Apesar disso, restou claro que a maioria esmagadora das respostas confluiu para as indicações do padrão apresentado, sendo relevante indicar que o tema remeteu para a luta em relação à igualdade de gênero, temática candente em especial para as concluintes do sexo feminino.
Diversos textos se apresentaram de forma concatenada, fundamentada e, muitas vezes, lamentando que, nos dias atuais, ainda se discuta a necessidade de proteger as mulheres dessa espécie de violência.
Em muitos casos, também, os concluintes fizeram um liame entre o acesso à educação e a melhora das relações entre homens e mulheres, ao permitir que a conclusão de maior grau de educação formal conduzisse à diminuição do índice de violência contra a mulher.
Foram identificadas respostas consoantes com o padrão apresentado, indicando os itens que nele constam em relação aos reflexos pessoais, na família e na sociedade.
As respostas dos estudantes raramente faziam referência à organização protetiva instituída pela Lei Maria da Penha, como esperado pelo padrão de resposta.
A respeito de tal aspecto o padrão indicava como uma das possibilidades de respostas ( item ‘ b ’ ) : a criação de casas de refúgio ou casa-abrigo, para acolher mulheres vítimas de violência doméstica ; reforço às Delegacias de Atendimento à Mulher ; inclusão da vítima em programas assistenciais do governo, programas de proteção à vítima e à testemunha, transferência de local de trabalho, se a vítima for servidora pública.
Como paradigma, indicaram -se as Delegacias da Mulher que são organizadas, com a quase integralidade dos seus membros de policiais do sexo feminino, e especializadas no atendimento a mulheres submetidas a situações de violência.
Essa preocupação gerencial tem por escopo bloquear os preconceitos que acompanham essas dramáticas histórias de violência, o que gera uma incompreensão e uma ineficiência na rede de proteção.
As respostas demonstram que os concluintes têm ciência do problema, da sua gravidade, do vínculo com a questão de gênero e com a submissão da mulher como efeito da sociedade patriarcal, ainda que tratado de forma superficial, com base na mídia e, algumas vezes, com apoio de experiências pessoais.
Identificou -se uma generalizada compreensão de que o problema está sendo enfrentado e que ainda existe um longo caminho no percurso para a sua resolução. 3.3.1.5 - Análise de Língua Portuguesa das Questões Discursivas do Componente de Formação Geral Os dados de Medicina, obtidos a partir das respostas às questões discursivas do Componente de Formação Geral, no que tange à Língua Portuguesa, encontram -se na Tabela 3.16 e no Gráfico 3.9.
Nesse aspecto, os alunos, de todo o Brasil, tiveram Média 68,3.
A maior Média com respeito à Língua Portuguesa foi obtida na região Sul ( 69,7 ), e a menor, na região Norte ( 65,9 ).
Quanto à variabilidade das notas, o Desvio padrão de todo o Brasil foi 17,5.
O menor Desvio padrão foi obtido na região Sul ( 16,5 ) e o maior Desvio padrão foi obtido na região Norte ( 19,0 ).
A Mediana das notas de Língua Portuguesa foi 72,5 para o Brasil como um todo.
A maior Mediana com respeito à Língua Portuguesa foi obtida na região Nordeste ( 75,0 ), e a menor, nas regiões Norte ( 70,0 ).
A nota Máxima para todo o Brasil foi de 95,0, com, pelo menos, um aluno tirando essa nota nas regiões Norte e Sudeste.
Além disso, a nota Mínima foi zero em todas as regiões do país.
O Gráfico 3.9 mostra a distribuição das notas de Língua Portuguesa do Componente de Formação Geral.
Observa -se que a maior frequência corresponde aos alunos que obtiveram nota no intervalo ( 70 ; 80 ].
Destacam -se, também, os alunos que deixaram ambas as questões em branco, representando quase 5 \% do total, caracterizando -se como máximo local. 3.3.1.6 - Comentários sobre a correção das respostas de Formação Geral com respeito à Língua Portuguesa Na avaliação do desempenho linguístico dos participantes foram consideradas quatro competências distintas, de modo a permitir um mapeamento detalhado do domínio dos recursos disponíveis na Língua Portuguesa para a comunicação escrita formal.
Nessa situação comunicativa, as exigências para o desenvolvimento desse tipo de texto dizem respeito : à adequação da seleção vocabular, ao desenvolvimento do conteúdo, à estruturação sintática dos períodos, à organização lógica das ideias, à utilização de procedimentos de encadeamento textual e de referenciação, à obediência às exigências morfossintáticas próprias da modalidade escrita da norma-padrão, ao respeito às regras ortográficas e às regras de acentuação gráfica.
A avaliação foi realizada considerando -se os textos das respostas às duas questões de Formação Geral e os resultados obtidos em cada uma delas, como seria de se esperar, foram semelhantes.
As competências avaliadas podem ser sintetizadas nos aspectos descritos a seguir.
Aspectos textuais - estruturação textual condizente com o gênero discursivo e o modo de organização textual expositivo/argumentativo adequado ao gênero – essa competência envolve : a estruturação sintática condizente com o padrão da modalidade escrita formal da Língua Portuguesa de modo a garantir a clareza necessária ; a distribuição do conteúdo do texto em parágrafos, de modo a garantir a sua organização temática ; a utilização de operadores discursivos que contribuam para a progressão temática do texto, estabelecendo relações lógicas entre as ideias apresentadas, tanto do ponto de vista intrafrasal, como do interfrasal ; a utilização de procedimentos de referenciação lexical e pronominal que permitam a retomada de referentes textuais ; o respeito às regras de pontuação como fator de estruturação do período.
Aspectos morfossintáticos - domínio dos diferentes aspectos morfossintáticos próprios da modalidade escrita formal da norma-padrão da Língua Portuguesa – essa competência envolve : a concordância nominal, a concordância verbal, a regência nominal, a regência verbal, a flexão nominal, a flexão verbal, a correlação entre os tempos verbais, a colocação pronominal e a utilização de sinais de pontuação que contribuam para a organização lógica da frase e do texto.
Aspectos vocabulares - seleção vocabular adequada à modalidade escrita formal da Língua Portuguesa, exigida pela situação comunicativa – essa competência envolve a precisão na seleção/utilização do vocabulário relacionado à temática solicitada pela questão ; a ausência de marcas da oralidade, como termos de sentido muito genérico ( “ coisa ”, “ negócio ”, “ você ” ) e termos de registros mais informais ( como gírias, jargões, frases feitas, ditados populares, termos regionais ).
Assim, espera -se que o participante respeite a adequação vocabular não usando gírias ou expressões coloquiais, evite repetição desnecessária de palavras e utilize um vocabulário mais formal, como solicitado por um texto dissertativo.
Aspectos ortográficos - respeito às convenções ortográficas da norma-padrão da Língua Portuguesa – essa competência envolve o domínio das regras de acentuação gráfica e da grafia padrão das palavras ( com ausência de abreviaturas próprias da linguagem da internet ), de acordo com as convenções estabelecidas pela legislação em vigor e consubstanciadas no Vocabulário Ortográfico da Língua Portuguesa, editado pela Academia Brasileira de Letras. ( Vale notar que, nessa edição de 2016, não será aceita a legislação anterior, no caso das regras relativas ao uso do hífen e da acentuação gráfica ).
A escolha dessas competências para subsidiar o processo de avaliação apoia -se na concepção de que, no desempenho dos graduandos, a modalidade escrita tem apresentado uma intensa simplificação, originada no padrão da modalidade oral da Língua Portuguesa.
No caso do texto de base dissertativa, inscrito em um registro formal, a distância entre as duas modalidades é ainda maior, o que provoca situações de hipercorreção ( desvios provocados pela incorporação indevida de uma regra da norma-padrão ) e de truncamentos sintáticos ( estruturas frasais incompreensíveis devido à complexidade sintática própria da modalidade escrita ).
Para efeito de pontuação, o critério aprovado pelo Inep, e já utilizado na avaliação do desempenho linguístico nas edições do Enade desde 2013, reúne as competências acima descritas em três grupos da seguinte forma : aspectos textuais ( 40 \% ), aspectos morfossintáticos e vocabulares ( 40 \% ) e aspectos ortográficos ( 20 \% ).
A seguir, analisam -se qualitativamente os resultados da avaliação de cada um dos aspectos avaliados nos textos de resposta das questões discursivas 1 e 2 de Formação Geral.
Aspectos textuais Esta competência é a que se revelou como a mais problemática entre os participantes, porque são muitos os desvios observados.
Em relação à estrutura formal do texto produzido pelos participantes, observaram -se relações linguísticas quase agramaticais, como as estabelecidas pela sequência de gerúndios sem o apoio de um ponto de partida para a organização das informações gramaticais e semânticas, ou seja, sem oração principal.
Observou -se que uma grande parte dos participantes não organizou as ideias em parágrafos, talvez devido ao pequeno número de linhas disponibilizadas para a resposta da questão ou, quem sabe, pela suposição de que não seria necessária essa divisão por não se tratar de um texto no modelo de uma redação dissertativo-argumentativa, como solicitado nos vestibulares.
Quanto à utilização dos mecanismos de referenciação, deve -se destacar a ocorrência, em uma parte dos textos, de repetições de palavras ou expressões sem a utilização de termos sinônimos ou pronomes, como seria adequado.
Outro aspecto relevante a ser destacado é a quase total ausência de operadores argumentativos, tanto intrafrasais como interfrasais, repercutindo uma tendência atual da mídia escrita.
A respeito do uso dos sinais de pontuação, observaram -se vários desvios nos textos analisados, com a ocorrência, inclusive, de parágrafos sem marca interna de pontuação para separar os períodos.
São os seguintes os tipos de problemas encontrados, que se caracterizam como desvios ao estabelecido pela norma-padrão : a ) uso inadequado de vírgula : • ocorrência entre o sujeito e o predicado, o que é considerado inadequado por interromper a sequência natural da frase ; • ocorrência de apenas uma das vírgulas para intercalar uma palavra, uma expressão ou uma oração encaixada, quando o correto seria a demarcação com uma vírgula antes e uma depois do termo encaixado na frase ; • ocorrência da vírgula no lugar do ponto para separar ideias que constituem períodos distintos ; b ) uso inadequado de ponto e vírgula para separar elementos que deveriam ser separados por vírgula, por estarem dentro do mesmo período ; c ) ausência de ponto final para separar períodos.
Vale observar que não foi penalizada a ausência de vírgula para destacar locuções ou adjuntos adverbiais de pequena extensão deslocados de posição na frase, por ser um uso opcional.
Os problemas relativos à estruturação textual geram, portanto, os seguintes tipos de desvios : • sequência justaposta de ideias sem encaixamentos sintáticos ; • redução de estruturas subordinadas, ao lado do aumento na frequência de estruturas coordenadas e absolutas ; • redução no uso de conectores para expressar relações lógicas essenciais à construção do texto ; • emprego inadequado de operadores que não estabelecem relações lógicas coerentes entre ideias do texto ; • repetição de termos sem a utilização de procedimentos mais sofisticados de substituição ( hiperonímias, hiponímias, nominalizações, expressões metafóricas ) ; • frases fragmentadas que comprometem a estrutura lógico-gramatical ; • frases formadas apenas por oração subordinada, sem oração principal.
Aspectos morfossintáticos e vocabulares Os resultados relativos aos aspectos morfossintáticos ( concordância, regência, colocação de pronomes, flexão nominal e verbal ) evidenciam desvios recorrentes no desempenho dos participantes, relacionados ao seu processo de formação e, principalmente, a hábitos da oralidade, marcados pela informalidade.
Apesar de exemplificarem processos de mudança em curso na Língua Portuguesa, esses desvios foram penalizados por não atenderem às exigências da norma-padrão relativas ao texto escrito em situação formal, como deveria ser o texto dissertativo solicitado no Enade.
O desvio mais frequente, em relação à regência, é a falta do sinal indicativo da crase – isso revela que o usuário não tem consciência de que, sob a forma do termo “ a ”, existe a presença de uma preposição “ a ”, exigida pela regência do termo anterior, combinada a um artigo definido ( por exemplo : “ Devido as grandes guerras que estão acontecendo... ” ) Outros problemas relacionados à regência verbal e à nominal foram identificados em alguns textos dos participantes : ( a ) a ausência de preposição ou seu uso indevido antes de pronome relativo ( por exemplo : “ a região que eles moram não oferece condições de segurança ” ) ; ( b ) o uso indevido do pronome relativo ( por exemplo : “ não pensam na cidade para que eles vão ” ) ; ( c ) utilização de pronome “ onde ” deslocado de sua função locativa ).
Esses processos refletem hábitos da modalidade oral da língua, em situações de registro informal.
Apesar da possibilidade de que essa alteração de regência se generalize no padrão escrito da Língua Portuguesa, como já está ocorrendo até em textos jornalísticos, a ausência da preposição foi penalizada neste processo de avaliação.
Outro desvio observado nos textos diz respeito aos processos de concordância verbal e de concordância nominal.
Quanto à concordância de número, além das ocorrências já conhecidas de falta da marca de plural em verbos ou adjetivos, em função da ocorrência de substantivo plural na função de sujeito, observaram -se dois tipos de desvio já registrados nas pesquisas linguísticas : ( a ) a ausência de marca de plural com sujeito posposto ( por exemplo : “ Foi expulso do país os imigrantes que entraram sem autorização da polícia ” ) ; ( b ) o uso indevido da marca de plural para estabelecer concordância com o plural da locução adjetiva, apesar de o substantivo que funciona como núcleo do sintagma nominal estar no singular ( por exemplo : “ A proteção das mulheres que sofrem violência foram garantidas pela Lei Maria da Penha ” ).
Uma ocorrência generalizada foi a ausência de acento circunflexo na forma plural do presente do indicativo do verbo “ ter ”, que foi considerada como um desvio na concordância verbal e não meramente de acentuação gráfica.
Quanto à concordância de gênero, alguns casos foram observados, normalmente no âmbito de sintagmas nominais longos, em que o adjetivo está afastado do substantivo.
Quanto à questão da colocação pronominal, foram poucos os desvios observados.
Apesar de serem aspectos relacionados à oralidade, observou -se que a maioria dos participantes já incorporou regras como a não introdução da frase por um pronome oblíquo e a próclise na presença de um termo atrator.
Não se adotou, entretanto, o padrão excessivamente formal descrito pelas gramáticas normativas em relação à posição do pronome oblíquo em locuções verbais, já que esse uso está muito distante da prática cotidiana, até em textos mais formais.
Quanto aos aspectos vocabulares ( precisão semântica do vocabulário utilizado e adequação ao registro formal exigido pelo tipo de texto solicitado no enunciado ), podem ser destacados os seguintes tipos de desvios : expressões da oralidade, seleção vocabular incompatível com o contexto, gerando situações de falta de inteligibilidade ; falta de domínio de vocabulário mais abstrato e de maior complexidade, essencial ao desenvolvimento do texto de base dissertativa.
Não se observou presença de palavras de baixa calão, ao contrário do exame de 2015, em que o tema “ funk ” gerou a utilização desse tipo de vocabulário.
Várias marcas de oralidade foram identificadas, embora não em alta frequência : o uso do pronome relativo “ onde ” como relativo universal, falta de artigo definido antes de substantivo, repetição de palavras por falta de vocabulário, reduções como “ tá ”, “ pra ”, “ pro ”, “ prum ”, expressões informais.
Aspectos ortográficos : O desempenho dos participantes revelou uma diferença muito grande nos dois aspectos de caráter ortográfico analisados neste item : grafia e acentuação.
Observou -se a ocorrência de baixo índice de desvios na grafia das palavras e grande índice de desvios no uso dos sinais de acentuação gráfica.
Em vários casos, ocorre ausência completa de acentuação.
Os resultados revelam que existe uma tendência crescente entre os universitários brasileiros de eliminação da acentuação gráfica, talvez motivada pelos hábitos relacionados às redes sociais ou devido a um conhecimento ainda superficial sobre o Acordo Ortográfico da Língua Portuguesa de 1990.
Os casos mais sistemáticos de eliminação do acento indicador da sílaba tônica são : • palavras proparoxítonas ( por exemplo : “ proximos ”, “ politicos ”, “ publica ”, “ fisica ”, “ psicologica ” ) ; • palavras paroxítonas terminadas em ditongo crescente ( por exemplo : “ necessario ”, “ noticiarios ”, “ individuo ”, “ dependencia ”, “ varios ”, “ propria ”, “ violencia ”, “ ocorrencia ”, “ suicidio ” ) ; • palavras oxítonas ( por exemplo : “ ate ”, “ tambem ”, “ mes ” ) ; • palavras paroxítonas terminadas em ditongo nasal ( por exemplo : “ orgão ”, “ orfão ” ).
Quanto ao domínio das convenções relativas à grafia das palavras, pode -se concluir que os desvios estão majoritariamente relacionados a questões relativas à oralidade : • hipercorreção pela escolha de “ e ” no lugar de “ i ” ou de “ i ” no lugar de “ e ”, por influência de hábitos da oralidade : por exemplo, “ inviado ” ( no lugar de “ enviado ” ), “ entervenção ” ( no lugar de “ intervenção ” ), “ empedimento ” ( no lugar de “ impedimento ) ; • eliminação do “ r ” marcador do infinitivo verbal : por exemplo, “ esta ” ( no lugar de “ estar ” ), “ estuda ” ( no lugar de “ estudar ” ), “ migra ” ( no lugar de “ migrar ” ) ; • utilização de “ s ” no lugar de “ c ” : por exemplo, “ mensionado ” ( no lugar de “ mencionado ” ) ; “ insentivou ” ( no lugar de “ incentivou ” ) ; • utilização de “ s ” no lugar de “ x ” : por exemplo, “ espectativa ” ( no lugar de “ expectativa ”, devido à confusão com o verbo “ esperar ” ) ; e a situação contrária, “ x ” no lugar de “ s ”, “ expantoso ” ( no lugar de “ espantoso ” ) ; • utilização de “ ss ” no lugar de “ ç ” : por exemplo, “ gerassões ” ( no lugar de “ gerações ” ) ; • uso de “ am ” no lugar de “ ão ” na forma de terceira pessoa do singular do presente e do pretérito perfeito do indicativo : por exemplo, “ estam ” ( no lugar de “ estão ” ) ; • utilização de “ l ” no lugar de “ u ” em final de sílaba : por exemplo, “ tralma ” ( no lugar de “ trauma ” ) ; • outros casos de desvios de grafia relacionados à variação diastrática podem ser observados em “ estrupo ” ( no lugar de “ estupro ” ), “ subjulgadas ” ( no lugar de “ subjugadas ” ), “ precoseito ” ( no lugar de “ preconceito ” ), “ soubre ” ( no lugar de “ soube ” ), “ vecendo ” ( no lugar de “ vencendo ” ), “ indesencia ” ( no lugar de “ \textbf{indecência} ” ), “ apolojia ” ( no lugar de “ apologia ” ), “ fulga ” ( no lugar de “ fuga ” ). • inadequação no uso da maiúscula para destacar determinadas palavras-chave do texto, como “ Violência ”, “ Brasileiros ”.
Destaque -se, também, a ocorrência de participantes que grafam os textos inteiramente em caixa alta.
Vale observar que, ao contrário do esperado, não apareceram abreviaturas próprias do “ internetês ”, relacionadas ao uso de redes sociais e de e-mails.
Sua influência, portanto, parece se restringir à ausência dos acentos gráficos.
Em relação aos aspectos textuais, observa -se dificuldade na estruturação textual, o que resulta em truncamentos sintáticos e prejuízo da compreensão do texto, e desrespeito a regras essenciais de pontuação, principalmente pela dificuldade em identificar as relações lógicas que se estabelecem entre as unidades frasais.
Em relação aos aspectos morfossintáticos, domínio das regras de concordância nominal e verbal, mas desrespeito às regras de regência nominal e verbal, principalmente no que diz respeito ao uso do sinal indicativo da crase.
Em relação aos aspectos vocabulares, reduzida utilização de palavras de uso mais formal e abstrato, adequado à construção do texto de base dissertativa.
Em relação às convenções ortográficas, domínio da grafia das palavras, mas falta de atendimento à maioria das regras de acentuação gráfica. 3.3.2 - Componente de Conhecimento Específico Na parte da prova relativa às questões discursivas do Componente de Conhecimento Específico ( Tabela 3.17 ), observa -se que a Média foi bem mais baixa do que para as questões discursivas do Componente de Formação Geral.
Enquanto no Componente de Formação Geral a Média para estudantes de Medicina de todo o Brasil foi 64,9, na parte de Conhecimento Específico a Média foi 32,0.
A maior Média deste componente foi obtida pelos estudantes da região Centro-Oeste ( 33,7 ), e a menor, pelos da região Norte ( 30,3 ).
Quanto à variabilidade das notas, o Desvio padrão de todo o Brasil foi 13,6.
O maior Desvio padrão foi encontrado na região Centro-Oeste ( 14,1 ), e o menor, na região Sul ( 13,1 ).
A maior nota Máxima, 78,3, foi obtida por pelo menos um estudante da região Sudeste.
Nas demais regiões a nota Máxima foi : 73,3 nas regiões Norte e Nordeste, 76,7 na região Sul e 75,0 na região Centro-Oeste.
Além disso, a nota Mínima ( 0,0 ) foi obtida por, pelo menos, um aluno em todas as regiões do Brasil.
A Mediana do Brasil como um todo foi 31,7.
A maior Mediana foi obtida nas regiões Nordeste, Sul e Centro-Oeste ( 33,3 ), e a menor, nas regiões Norte e Sudeste ( 30,0 ).
O Gráfico 3.10 representa a distribuição das notas nas questões discursivas no Componente de Conhecimento Específico.
A moda desta distribuição ocorre no intervalo ( 20 ; 30 ].
Destacam -se, também, os alunos que obtiveram nota no intervalo [ 30 ; 40 ] com aproximadamente 26 \% do total.
A análise de cada uma destas questões será feita a seguir. 3.3.2.1 - Análise da Questão Discursiva 3 do Componente de Conhecimento Específico Na questão 3, cujos resultados aferidos encontram -se descritos na Tabela 3.18, a Média dos estudantes de todo o Brasil foi 18,3.
A menor Média nessa questão foi obtida pelos alunos da região Norte ( 16,5 ), enquanto a maior Média foi obtida na região Centro-Oeste ( 21,3 ).
Quanto à variabilidade das notas, o Desvio padrão de todo o Brasil foi 17,8.
O maior Desvio padrão foi obtido na região Centro-Oeste ( 18,6 ), enquanto o menor foi obtido na região Norte ( 17,3 ).
A nota Máxima, 80,0 pontos, foi alcançada por, pelo menos, um aluno de todas as regiões, com exceção da região Norte, onde a nota Máxima foi 75,0.
A Mediana do Brasil como um todo foi 10,0, a mesma para todas as regiões.
A nota Mínima ( 0,0 ) foi obtida por, pelo menos, um estudante em todas as regiões do Brasil.
O Gráfico 3.11 mostra a distribuição das notas na questão discursiva 3, do Componente de Conhecimento Específico da área de Medicina.
Poucos estudantes deixaram a questão em branco, mas muitos receberam nota zero, quase a mesma quantidade daqueles que obtiveram notas no intervalo modal ( 0 ; 10 ].
Destaca -se, ainda, outro máximo local no intervalo ( 30 ; 40 ]. 3.3.2.2 - Comentários sobre as respostas à Questão Discursiva 3 A Questão 3 abordou tema relativo ao diagnóstico e à prevenção de patologias cujo conteúdo enquadra -se nas Diretrizes Curriculares para o curso de graduação em Medicina.
No caso em foco, o diagnóstico precoce é fundamental para o correto encaminhamento do atendimento das crianças e de suas famílias, com possibilidade de melhorar o prognóstico dos recém-natos.
A recente epidemia de microcefalia, associada à infecção por Zika vírus, colocou essa questão em grande evidência.
Porém, mesmo na ausência de situações excepcionais, é cada vez maior a necessidade de atenção e competência no diagnóstico precoce de anomalias congênitas.
Com a importante queda, nos últimos trinta anos, da mortalidade infantil por causas passíveis de prevenção no Brasil, através de medidas de saúde pública, como vacinação e saneamento básico, cresce, em importância, o percentual da mortalidade infantil atribuível a doenças e anomalias congênitas.
Assim, é extremamente relevante a competência dos médicos e outros profissionais de saúde no diagnóstico e correto encaminhamento dessas patologias.
A questão tem enunciado claro e solicita a descrição de procedimento técnico padrão, que deve ser aplicado rotineiramente no atendimento aos recém-natos, incluindo sua interpretação e o encaminhamento dos neonatos que forem diagnosticados com possível comprometimento neurológico, resultante de infecção pelo vírus Zika.
Apesar de a relevância deste procedimento estar em especial evidência no momento, como consequência da emergência sanitária decretada em função da epidemia por Zika vírus, o procedimento integra as ações previstas no exame de rotina do recém-nato em qualquer circunstância.
Além dos procedimentos que orientam o diagnóstico inicial ( itens ‘ a ’ e ‘ b ’ ), é exigido que o respondente explique as etapas do atendimento previsto para essas crianças pelo Ministério da Saúde.
Os documentos oficiais que tratam dessas ações tiveram ampla divulgação, especialmente voltada para os profissionais de saúde e as escolas médicas desde 2015.
Os conhecimentos exigidos, tanto teóricos como de técnicas, e o nível de profundidade previsto pelo padrão de resposta são compatíveis com o esperado ao final do curso de formação de médicos.
A questão exigiu a descrição de uma técnica simples de exame do recém-nato, a interpretação de seus resultados, e o encaminhamento que deveria ser feito para o correto acompanhamento dos neonatos com patologia neurológica.
O item ‘ a ’ exigiu a descrição simples da técnica de medição do perímetro cefálico.
Não há possibilidade de incompreensão do enunciado.
Alguns estudantes apresentaram dificuldade de oferecimento de técnica, mas, grande parte ofereceu respostas compatíveis com o padrão esperado.
O item ‘ b ’ também não mostrou maior dificuldade para os estudantes em relação à compreensão do enunciado, mas mostrou uma grande diversidade em relação à escolha de padrões de comparação para definição da microcefalia.
As respostas ao item ‘ c ’ também não indicaram incompreensão do enunciado e, sim, a falta de conhecimento em relação ao correto encaminhamento dos neonatos com comprometimento neurológico.
Tratando -se de tema, atualmente, de ampla divulgação, inclusive pela imprensa leiga, grande número de estudantes parece não ter familiaridade com esse tema, que vem sendo amplamente divulgado pelo Ministério da Saúde desde 2015.
Em suma, como o exigido nessa questão é o conhecimento de técnicas e padrões, não há muito espaço para a diversidade de respostas, e o observado foi uma grande uniformidade das respostas corretas e, também, das incorretas.
As respostas incorretas do item ‘ a ’ mostram desconhecimento dos acidentes anatômicos que devem ser considerados na medida do perímetro cefálico.
As do item ‘ b ’ mostram o desconhecimento dos padrões internacionalmente adotados para a definição de microcefalia.
As do item ‘ c ’ mostram o desconhecimento das normas estabelecidas no Brasil pelo Ministério da Saúde para o acompanhamento das crianças com suspeita de microcefalia.
A linguagem adotada pelos alunos nessa questão foi, em geral, adequada quanto ao uso de termos específicos da área.
No entanto, nessa, como nas demais questões, observamos um grande número de respostas de difícil leitura, seja por problemas de redação ou de ' letras incompreensíveis ' ; mesmo com o uso dos bons recursos de aumento de imagem e de maior realce do texto que estavam disponíveis para os corretores.
Este tem sido um problema recorrente na avaliação de questões com respostas discursivas manuscritas.
Considero importante enfatizar esse ponto em relação à formação médica, pois, ainda hoje, esta é a principal forma de comunicação do profissional médico com farmácias, laboratórios de análises clínicas e outros profissionais.
Nessas circunstâncias, a eventual incompreensão do texto pode gerar situações de risco para os pacientes.
Nessa questão, as respostas consideradas fracas mostram a falta de familiaridade com o procedimento de medição do perímetro cefálico, indicando a falta de exposição a este treinamento durante o curso, pois se trata de procedimento que, uma vez realizado, dificilmente seria esquecido.
Todos os itens dessa questão se referem a conhecimentos que exigem um treinamento em serviço.
A observação destes procedimentos e condutas torna fácil a sua memorização.
Por outro lado, a memorização na ausência da observação me parece extremamente difícil.
As respostas medianas mostram um índice de acertos um pouco maior, mas não diferem significativamente das respostas fracas.
Em geral, vemos, neste grupo, respostas melhores em relação aos itens ‘ a ’ e ‘ b ’ quando comparamos com as respostas fracas.
A maioria das respostas consideradas boas mostrou melhor desenvolvimento do item 3 - normas de acompanhamento da criança com suspeita de microcefalia.
No entanto, a maioria dos estudantes teve desempenho fraco nessa questão.
Não houve grande diversidade de abordagem em relação a essa questão.
Praticamente, todos os alunos pareceram ter entendido o que estava sendo solicitado e descreveram a medida do perímetro cefálico, os critérios estatísticos de definição de microcefalia atualmente adotados ( tabelas da OMS e MS ), e as orientações que deveriam ser observadas para o correto encaminhamento das crianças, após o diagnóstico inicial.
No entanto, analisando as respostas, ocorreram alguns erros em relação à descrição dos pontos cranianos, que servem de referência para a medição do perímetro.
Mais do que desconhecimento da ' nomina anatomica ', as respostas indicaram certa deficiência dos respondentes quanta à experiência prática nesse campo.
E ainda, em algumas respostas, houve certa confusão entre os conceitos de ' perímetro ' e ' diâmetro '.
As respostas relativas aos padrões de definição de microcefalia revelaram uma variedade de padrões de referência, ainda que a maioria soubesse que era preciso comparar a medida obtida com um padrão determinado.
Quanto ao encaminhamento das crianças, foi notado uma diversidade grande de propostas de acompanhamento dos casos, sugerindo que este ponto poderia não ter sido abordado de forma estruturada durante o curso.
Com a progressiva diminuição dos índices de mortalidade infantil no Brasil, como em outros países, é também observada uma mudança no perfil de causas de morte, e as anomalias congênitas tendem a aumentar sua participação na mortalidade proporcional.
Isto não significa, necessariamente, uma maior incidência desses eventos.
Assim, mais do que abordar um tema atual, o conhecimento avaliado nessa questão se refere aos conteúdos e às competências indispensáveis para os profissionais que estarão em atividade nos próximos anos.
O acompanhamento da microcefalia era assunto tratado por especialistas e ocupava espaço restrito na formação profissional.
No entanto, seu diagnóstico inicial, a partir da medida do perímetro cefálico, sempre teve importância na formação e deve ser abordado, pelo menos em duas disciplinas do curso médico – a obstetrícia e a pediatria.
Os erros cometidos na descrição da medida do perímetro cefálico sugeriram que a maioria dos alunos nunca realizou este procedimento e, provavelmente, poucas vezes o viu ser realizado.
Da mesma forma, a falta de informação sobre como encaminhar os casos parece indicar uma falha do treinamento na assistência pediátrica na atenção básica.
O conteúdo das respostas, fortemente, indica a necessidade de avaliação e revisão desses aspectos da formação.
As atuais Diretrizes Curriculares enfatizam o treinamento em cenários de prática, especialmente na atenção básica.
Essa questão avalia conhecimentos para cuja fixação é essencial este treinamento. 3.3.2.3 - Análise da Questão Discursiva 4 do Componente de Conhecimento Específico A Tabela 3.19 contém as informações relativas à questão 4 do conjunto de questões do Componente de Conhecimento Específico.
O desempenho dos estudantes nessa questão foi superior ao da questão 3 ( Média 18,3 ) e inferior ao da questão 5 ( Média 50,8 ).
A Média geral do Brasil foi 26,9, sendo a menor Média registrada na região Sudeste ( 26,1 ), e a maior, nas regiões Norte e Nordeste ( 28,3 ).
A nota Máxima ( 95,0 ) foi atingida por, pelo menos, um concluinte da região Sudeste.
A Mediana em todo o Brasil foi 25,0, o mesmo valor foi obtido nas regiões Sudeste, Sul e Centro-Oeste.
Nas regiões Norte e Nordeste a Mediana foi maior ( 30,0 ).
As notas mínimas foram zero em todas as regiões.
O Gráfico 3.12 representa a distribuição de notas da questão discursiva 4, do Componente de Conhecimento Específico.
Essa distribuição é unimodal com moda no intervalo ( 10 ; 20 ]. 3.3.2.4 - Comentários sobre as respostas à Questão Discursiva 4 A Questão Discursiva 4 aborda tema fundamental para a formação médica, especialmente para a prática em centros urbanos – a atenção imediata, emergencial, a vítimas de acidentes de trânsito.
A atuação na atenção a emergências integra as diversas competências em relação aos procedimentos de diagnóstico e de tratamento previstos nas Diretrizes Curriculares, sendo previsto que “ o mínimo de 30 \% ( trinta por cento ) da carga horária prevista para o internato médico da Graduação em Medicina será desenvolvido na Atenção Básica e em Serviço de Urgência e Emergência do SUS, respeitando -se o mínimo de dois anos deste internato. ” Os conhecimentos exigidos pelo enunciado da questão integram habilidades de realização de anamnese, exame físico, diagnóstico e intervenção, num cenário que exige rapidez de resposta e priorização das intervenções essenciais à manutenção da vida e à melhoria do prognóstico do paciente.
Para o treinamento dos alunos e para a prática nos cenários de emergência, foram desenvolvidos protocolos de atuação que são adotados internacionalmente e que devem ser conhecidos pelos alunos.
Este conhecimento está sendo verificado nessa questão.
O caso apresentado descreve um acidente de trânsito em que a vítima é um homem de 25 anos, levado de ambulância ao serviço de emergência de um hospital.
Como é comum nesses casos, a anamnese é limitada à descrição das circunstâncias do acidente pelos paramédicos que fizeram o primeiro atendimento e transportaram o acidentado até o hospital.
Os conteúdos aqui avaliados são os de competência diagnóstica rápida e, a partir de exame sumário do paciente, definição de intervenção emergencial e execução de procedimentos compatíveis com o diagnóstico imediato.
Definir rapidamente as ações e seguir uma sequência definida de “ passos do atendimento ” é essencial para aumentar a probabilidade de sobrevivência do paciente.
Estes conteúdos integrados em protocolos de atendimento a emergências são hoje amplamente utilizados no treinamento profissional nas escolas médicas.
O nível de profundidade esperado pelo padrão de resposta exige estes conhecimentos, mencionando ( nomeando ) ou não o protocolo que está sendo seguido.
O enunciado da questão é claro, e as informações ali contidas permitem responder adequadamente ao que é solicitado.
Em geral, os estudantes parecem ter compreendido plenamente o que era solicitado e não registramos interpretações diferentes do esperado pelo padrão adotado.
Ressaltamos que essa questão – que cobra a descrição de uma sequência de procedimentos que é rigidamente padronizada – praticamente não permite interpretações diferentes da esperada.
Essa questão, que considerávamos mais difícil que a Questão 3 – ainda que mais fácil que a Questão 5 –, pelas notas obtidas, se mostrou mais fácil que a anterior para os alunos.
Provavelmente, a cobrança de um protocolo de intervenção amplamente conhecido seja a razão do melhor desempenho quando essa questão é comparada com a Questão 3.
É possível que, também, interfira no melhor desempenho ( ou, pelo menos, na tentativa ) a proximidade maior desse cenário de prática com o que é visto como mais representativo da “ ação do médico ”, já que este é, por exemplo, o cenário preferencial de séries de televisão de ampla audiência.
A maioria dos alunos conseguiu nomear o protocolo ( ABCDE ) que deve ser seguido no atendimento.
Apesar disso, a descrição correta dos passos previstos no protocolo foi feita pela minoria dos alunos.
Nesta, como na Questão 3, as respostas pareciam indicar pouca experiência direta com a situação apresentada, refletindo, nas respostas melhores, um conhecimento teórico do protocolo de atendimento na sala de emergência.
Não identificamos respostas em que o aluno tenha acrescentado elementos voltados à sua experiência pessoal em situações semelhantes.
Essa observação é especialmente preocupante porque o atendimento a emergências é, destacadamente, um desempenho que envolve habilidades cuja proficiência depende criticamente do treinamento em serviço.
O padrão de resposta prevê a descrição das etapas de um protocolo de atendimento conhecido como o “ ABCDE do trauma ” que descreve uma sequência de passos a serem seguidos no atendimento de emergência.
Essa sequência prioriza ações que garantem ao paciente a melhor chance de sobrevida com a menor incidência de sequelas, e as letras representam as etapas do atendimento nomeadas em inglês – A ( airways ) B ( breathing ) C ( circulation ) D ( disability ) E ( exposure ).
As iniciais em ordem alfabética têm o propósito de enfatizar a necessidade de aplicação das ações nessa sequência.
A maioria dos alunos mostrou ter conhecimento do protocolo, citando o “ ABCDE do trauma ”.
No entanto, poucos alunos conseguiram descrever corretamente as etapas.
A ordem das etapas também não parece ser um elemento percebido como essencial pela maioria.
Nenhum aluno conseguiu descrever corretamente todas as etapas na ordem prevista.
Como se pode ver pela distribuição de notas, poucos conseguiram descrever, adequadamente, a maioria das etapas.
As notas refletem menos os erros do que omissões importantes na descrição das ações previstas nas etapas de atendimento.
Na situação apresentada nessa questão, a omissão ( ou a mudança de ordem ) de uma etapa pode significar o óbito do paciente.
Aqui, de forma mais evidente ainda do que na Questão 3, parece faltar a exposição dos alunos, pelo menos a observação de situações como a descrita no enunciado.
Em poucas respostas, conseguimos perceber que o aluno já esteve em contato, como profissional em treinamento, com a situação descrita.
A linguagem dos alunos é adequada quanto ao uso de termos médicos específicos.
No entanto, a redação é, frequentemente, pouco clara, e a escrita é de muito difícil leitura.
Aparentemente, continua sendo comum a “ letra de médico ”, apesar de ser infração prevista no código de ética profissional a falta de clareza e legibilidade em documentos como receitas e atestados.
O desempenho dos alunos foi, em geral, fraco, ainda que bem superior ao observado em relação à Questão 3.
Apesar de cerca de 60 \% das questões mostrarem fraco desempenho, outros 38 \% dos alunos tiveram desempenho mediano e apenas 2 \% mostraram bom desempenho.
As respostas consideradas fracas se caracterizam pela descrição fora de ordem dos procedimentos previstos no protocolo ou com descrição insuficiente e inadequada como “ fazer palpação ” ou “ procurar déficit neurológico ” sem desenvolvimento adequado do procedimento.
No item ‘ b ’, foi grande o não entendimento do que era solicitado – o diagnóstico diferencial.
E, muito frequente, o erro de interpretação do som provocado pela percussão.
Este foi um resultado totalmente inesperado, já que se trata de procedimento simples e que integra ( ou deveria integrar ) o ensino da anamnese e exame físico padrão no curso médico.
Na resposta ao item ‘ c ’, poucos alunos acertaram, integralmente, a descrição das etapas do procedimento de drenagem torácica.
A assepsia e a anestesia local foram as etapas mais frequentemente omitidas.
As respostas medianas e boas mostram, progressivamente, maior organização na descrição dos procedimentos, especialmente quanto ao detalhamento do exame físico e do uso de parâmetros de avaliação do nível de consciência do paciente ( escala de Glasgow ).
As respostas boas acrescentam a esses aspectos uma descrição mais adequada do procedimento de toracocentese.
A abordagem mais frequente foi citar, inicialmente, o “ ABCDE do trauma ” e, em seguida, tentar descrever as etapas.
Essa pergunta não permite muita escolha quanto ao desenvolvimento da resposta, já que cobra o conhecimento de um protocolo, e isto foi identificado pela maioria dos alunos.
Os principais erros foram de omissão da descrição de procedimentos indispensáveis.
A descrição errada de procedimentos foi menos frequente, mas também ocorreu de forma significativa, especialmente quanto à descrição do procedimento cirúrgico ( toracocentese ).
A descrição equivocada da sequência de procedimentos também foi observada com frequência.
As respostas analisadas demonstraram frequente omissão quanto à descrição das etapas de assepsia e anestesia em um procedimento cirúrgico realizado em ambiente hospitalar.
Como se trata de procedimentos que exigem treinamento em serviço para que se possa atingir competência mínima, o teor das respostas indica que este treinamento não está sendo adequado para a maioria dos alunos.
Em relação ao conteúdo avaliado nessa questão, é importante observar que, dificilmente, apenas o estudo teórico levará ao bom desempenho.
É fundamental que o aluno domine os conhecimentos de anatomia, fisiologia e clínica médica, necessários ao bom desempenho neste cenário.
Mas, somente o treinamento em serviço é capaz de produzir a competência essencial para tomar as decisões e realizar os procedimentos – que integram aqueles conhecimentos – cuja descrição é solicitada nessa questão.
Assim, os resultados dessa avaliação indicam a necessidade de revisão do ensino desses conteúdos, especialmente quanto aos cenários de prática.
É importante enfatizar que a situação apresentada nessa questão é o dia a dia das emergências nas áreas urbanas do nosso país.
Também relevante apontar a aparente falta de competência ( ou, pelo menos, o esquecimento ) de grande parte dos alunos em relação a um procedimento simples do exame físico – a percussão.
A excessiva dependência de exames complementares poderia estar levando ao abandono de uma técnica simples e altamente eficiente em situações como a aqui descrita.
A percussão é executada pelo médico imediatamente no curso do atendimento, levando a uma intervenção imediata e bem orientada.
Com frequência, o exame de imagem não está tão rapidamente disponível, e a demora deste diagnóstico e a consequente demora da intervenção podem ter importante consequência sobre a letalidade na situação descrita. 3.3.2.5 - Análise da Questão Discursiva 5 do Componente de Conhecimento Específico A Tabela 3.20 contém as informações relativas à Questão 5 do conjunto do Componente de Conhecimento Específico.
O desempenho dos estudantes de todo o Brasil nessa questão foi mais alto dentre as três questões discursivas desse componente.
A nota Média dos estudantes de todo o Brasil foi 50,8.
A maior Média foi registrada na região Centro-Oeste ( 52,7 ), enquanto a menor Média foi registrada na região Norte ( 46,1 ).
Quanto à variabilidade das notas, o Desvio padrão dos alunos do Brasil, como um todo, foi 20,9.
Enquanto o maior desvio foi encontrado na região Nordeste ( 21,8 ), o menor foi encontrado na região Sul ( 20,4 ).
A Mediana para o Brasil foi 50,0, a mesma obtida em todas as regiões.
Para o conjunto de alunos de Medicina do Brasil, a nota Máxima foi 100,0, obtida por pelo menos um estudante das regiões Nordeste e Sudeste.
A nota Mínima ( 0,0 ) foi a mesma em todas as regiões.
O Gráfico 3.13 representa a distribuição das notas da questão discursiva 5 do Componente de Conhecimento Específico.
Essa distribuição tem moda no intervalo ( 40 ; 50 ].
O gráfico mostra, também, que nessa questão as notas ficaram mais espalhadas do que nas duas anteriores. 3.3.2.6 - Comentários sobre as respostas à Questão Discursiva 5 Os conhecimentos exigidos pelo enunciado da questão e pelo padrão de resposta estão totalmente compatíveis com os conteúdos das Diretrizes Curriculares da Medicina, tendo nível de profundidade adequado ao que se espera de um recém-formado.
Essa questão discursiva apresentou ao aluno uma situação que, provavelmente, representa a melhor aproximação a situações reais, que serão enfrentadas pela maioria dos formandos em medicina nos primeiros anos de prática.
O enunciado descreve um paciente do sexo masculino, de 52 anos, fumante, hipertenso, apresentando um quadro de insuficiência cardíaca.
A hipertensão arterial atinge, hoje, quase 1/4 da população adulta no Brasil, e o fumo contribui, de modo relevante, para complicações cardiovasculares nesses pacientes.
A mortalidade por doenças cardiovasculares atinge quase 30 \% dos óbitos no Brasil, sendo a principal causa de morte no país.
O progressivo aumento do número relativo de idosos na população do Brasil ( e mundial ) e a manutenção dos altos índices de doença cardiovascular sinalizam que essa questão é oportuna e representa uma situação comum da prática profissional.
O paciente descrito procura atendimento ambulatorial.
O enunciado da questão descreve aspectos da anamnese e do exame físico suficientes para a formulação de hipótese inicial de diagnóstico e para estabelecimento de conduta terapêutica imediata.
Os conhecimentos exigidos para tal integram as Diretrizes Curriculares da Medicina e o currículo dos cursos de medicina no Brasil.
O nível de profundidade dos conhecimentos necessários para responder ao que é solicitado é totalmente compatível com o conhecimento esperado de um graduando do curso de medicina.
O comando da questão requer a descrição de quatro medidas terapêuticas não farmacológicas e três medidas terapêuticas farmacológicas indicadas para este paciente.
Entre as medidas não farmacológicas listadas no padrão de resposta está a vacinação contra influenza e pneumococo.
Os raros casos de presença da indicação de vacinas para esse paciente seguiram o previsto no padrão de resposta – ou seja, as vacinas foram indicadas como terapêutica não farmacológica.
Essa confusão, certamente, reflete em questões que não são de rápida e simples solução, mas que, especialmente em relação às doenças não infecciosas, merecem um aprofundamento da discussão.
Assim, estamos acostumados a considerar, para prevenção das doenças não infecciosas ( ou para limitação de seu agravamento ), basicamente, as ações de “ mudança do estilo de vida ” que dependem de ações individuais ou coletivas referidas em grande parte à dieta – no sentido hipocrático do termo.
Ou seja, mudanças na alimentação, desempenho físico, e outros comportamentos que integram o conceito de “ estilo de vida. ” Já para as doenças infecciosas, consideramos, basicamente, intervenções farmacológicas.
Mesmo quando se consideram intervenções no meio ambiente, estas costumam ser pontuais, como, por exemplo, eliminar os mosquitos transmissores de dengue, febre amarela, etc.
E essas medidas são vistas como muito menos efetivas que as farmacológicas – as vacinas.
Para o tratamento, nos dois casos, vemos principalmente uma intervenção farmacológica.
Ainda que, evidentemente, hoje o tratamento ( e não apenas a prevenção ) da hipertensão e do diabetes ( as principais doenças não infecciosas ) utilize amplamente o recondicionamento físico através de exercícios programados e supervisionados por médicos especialistas – ou seja, um tratamento não farmacológico.
No mais, o enunciado é claro e fornece informações e dados suficientes para as respostas esperadas.
Pelas notas obtidas, é aparente que a maioria estava mais preparada para responder a essa questão do que às duas anteriores.
É interessante observar que essa questão avalia conhecimentos mais relacionados com a prática clínica, indicando que esta continua sendo o foco central do ensino da medicina.
Consideramos que isto é um ponto positivo.
Mas, ainda que essa questão mostre melhor desempenho que nas respostas às anteriores, a média das notas foi muito baixa e aponta para a necessidade de ampla reflexão sobre o padrão de formação dos médicos.
Em geral, os alunos elaboraram respostas que correspondem ao padrão esperado.
Entretanto, em relação às medidas não farmacológicas, a maioria citou a dieta hipossódica e o abandono do fumo.
Poucos mencionaram a reabilitação cardíaca e, mais rara ainda, foi a menção de vacinação contra influenza e pneumococo.
Ou seja, não houve grande diversidade de respostas e, sim, um número grande de respostas incompletas nesse item.
Em relação a medidas farmacológicas, a diversidade foi bem maior, especialmente quanto às explicações das indicações terapêuticas.
A maioria menciona, corretamente, pelo menos uma droga, mas, na explicação do uso, se limita a dizer que “ aumenta a sobrevida ”.
A explicação correta do mecanismo de ação como justificativa do emprego da droga – como previsto no padrão de resposta – é encontrada na minoria das provas.
Alguns alunos indicam drogas e dão explicações de uso que podemos classificar como “ do senso comum ”, como a indicação de uso de AAS com a justificativa de que “ afina o sangue ”.
Observamos, também, uma incidência inesperada de propostas terapêuticas em relação à presença de edema em membros inferiores e de bolsa escrotal.
Em relação a este último sinal, muitas provas apontam o uso de suportes compressivos e, para o edema de membros inferiores, a adoção de posição de Trendelenburg.
Em nenhum momento, essas respostas parecem formas de evitar responder adequadamente ao que é solicitado.
As explicações farmacológicas corretas para a indicação dos medicamentos são dadas pela minoria dos alunos.
Não se observou grande diversidade de padrão de resposta.
O que se observou foi que, mesmo fazendo a correta indicação terapêutica, a maioria desconhece o mecanismo de ação das drogas receitadas.
Muitas vezes, o uso de um diurético é justificado simplesmente “ para aumentar a diurese ”.
O de um betabloqueador para “ aumentar a sobrevida ”.
Também foi observado um “ excesso de resposta ”, ou seja, a indicação de mais do que as três medidas solicitadas.
Neste caso, o procedimento adotado foi considerar primeiro as corretas e ignorar as demais.
Em nenhum caso, identificamos incompreensão do enunciado dessa questão.
A linguagem dos alunos revela, em muitos casos, um processo de memorização, sem compreensão dos processos fisiológicos e farmacológicos envolvidos na prescrição de medicamentos.
E, também, pouca leitura.
É difícil entender que um aluno que indica um diurético “ polpador ” de potássio tenha estudado o suficiente sobre este medicamento.
No que diz respeito aos termos específicos da área, a maioria dos alunos parece estar, apenas, começando a se familiarizar com os mais comuns.
Apesar dos problemas já descritos, essa é a questão que mostrou melhores resultados, pois, apenas cerca de 20 \% dos alunos tiveram notas consideradas fracas.
Mais de 60 \% conseguiram desempenho mediano e, em torno de 15 \%, desempenho bom.
De fato, essa questão, que inicialmente nos pareceu apresentar maior dificuldade para os alunos, mostrou os melhores resultados.
É possível que isto reflita não apenas a relevância dada à clínica médica de adultos na formação profissional, como, também, a experiência dos alunos.
É provável que, mais de uma vez, durante o treinamento nos ambulatórios e enfermarias, os alunos tenham encontrado pacientes semelhantes ao descrito nessa questão.
É mesmo quase impossível que isto não tenha acontecido.
Assim, apesar da maior complexidade deste caso em relação aos apresentados nas outras duas questões, a experiência dos alunos e, consequentemente, o maior tempo de estudo dedicado aos aspectos clínicos de casos semelhantes podem compor a principal razão dessa diferença.
Apesar disso, em relação ao tratamento, mesmo em respostas boas, observamos indicações que apontam para falhas importantes nos conhecimentos de farmacologia, especialmente quanto à indicação simultânea de vários medicamentos.
Nessa questão, observamos, com maior nitidez, uma “ gradação de respostas ”.
Ou seja, as respostas fracas apontam para pouco ou nenhum conhecimento em relação à abordagem diagnóstica e terapêutica deste paciente.
As respostas medianas indicam, em geral, melhor conhecimento da abordagem diagnóstica do que das medidas terapêuticas.
As respostas boas indicam um domínio bom da abordagem diagnóstica e terapêutica em relação ao caso descrito.
Aparentemente, a abordagem mais frequente para responder a essa questão foi a de identificação de sinais e sintomas e de descrição de propostas de intervenção em relação aos sinais e sintomas descritos.
Poucos alunos, mesmo nas melhores respostas, apontaram, claramente, uma hipótese diagnóstica.
A maioria identificou sinais e sintomas em relação aos quais propôs intervenções terapêuticas.
Nas melhores respostas, encontramos descrições corretas dos efeitos esperados a partir das intervenções.
Das boas para as piores respostas, observamos um “ gradiente de acertos e erros ”, mais do que uma diversidade de propostas diagnósticas e terapêuticas.
Os principais erros parecem apontar para falhas de conhecimento da fisiopatologia – como a sugestão de elevação dos membros inferiores para tratamento do edema – e, também, da farmacologia – com a prescrição simultânea de medicamentos, sem justificativa adequada.
As respostas a essa questão parecem indicar, quando comparadas às respostas das anteriores, que há uma maior ênfase do treinamento clínico ( especialmente da clínica de adultos ) ambulatorial ou em internação, em relação ao treinamento de atenção às emergências ou na atenção básica.
Apesar disso, é bastante evidente a deficiência em relação a conhecimentos de farmacologia e fisiologia, gerando propostas terapêuticas conflitantes entre si ou que podem ter como resultado o agravamento do quadro do paciente.
A ênfase maior no treinamento clínico deveria se refletir num maior domínio da fisiopatologia.
No entanto, as propostas terapêuticas e suas justificativas não refletem isso, pelo contrário, indicam falta de treinamento clínico adequado, apesar de a prova indicar que, nessa área, ainda encontramos os melhores desempenhos dos alunos.
Apesar de várias questões objetivas abordarem situações, em parte, semelhantes às dessa questão, não identificamos nenhuma situação em que isto gerasse uma relação de influência em relação às respostas. 3.3.3 - Considerações Finais Pelo exposto, nota -se a necessidade de disponibilidade de estágio com atividades programadas e supervisão adequada, e, não, os inúmeros “ plantões ”, em unidades de saúde públicas ou privadas, que, sem adequada supervisão, não contribuem para melhor formação profissional, com isso, haveria um reflexo positivo em todo o sistema de atenção à saúde e sua disponibilidade para o treinamento profissional.

% matched lemmas: aportar, autoestima, homicídio, humilhação, indecência
\end{textsample}
